% LOCKED SCAFFOLD - Volume I FINAL TOC - Chapter 1 only
% Do not add Ch 2+ content to this file
\chapter{Foundations of Angular Electromagnetic Structure}
\label{ch:01}
\setcounter{section}{-1}

\section{Introduction and Scope}
\label{sec:ch01.0}

Chapter~\ref{ch:01} opens the monograph by preparing a rigorous reading frame
for all later derivations. The immediate task is not to enumerate modal families
or applications, but to separate three layers that are frequently mixed in
electromagnetic writing: governing field law, representation choice, and
operationally admissible observable content. Subsection~\ref{sec:ch01.0.1}
therefore states the chapter mandate and scope boundary,
Subsection~\ref{sec:ch01.0.2} develops angular structure as a consequence of
Maxwell dynamics under geometry and boundary constraints, and
Subsection~\ref{sec:ch01.0.3} defines the admissibility domain in which later
claims remain mathematically and physically meaningful.


\subsection{Scope and role of Chapter 1}
\label{sec:ch01.0.1}
% TOC tag: [HOME]

This chapter serves advanced graduate readers and researchers who already command
core Maxwell theory, yet require a proof-structured account in which each claim
is explicitly tied to assumptions and inheritance logic
\cite{jackson1999,harrington2001,collin2001}. Inside Volume~I, its role is
architectural: it creates the inferential vocabulary used in every later chapter.

That vocabulary rests on four commitments. Field laws are read in one fixed sign
convention and SI scaling; transport arguments remain conservation-governed;
representation-level language is kept subordinate to invariant content; and
admissibility is treated as both a mathematical and physical requirement.
Without this discipline, downstream operator and spectral statements would become
ambiguous in interpretation even when formally correct.

The scope boundary is equally strict. Chapter~\ref{ch:01} does not attempt the
operator-domain machinery developed in Chapter~2, the angular-eigenfield
construction of Chapter~3, or the finite-geometry accessibility limits of
Chapter~4. Those results remain at their dedicated owner locations and are inherited
forward by reference only. This partition is part of the no-recursion design and
is therefore methodological, not stylistic.

The subsection delivers one concrete outcome: all later text in
Chapter~\ref{ch:01} is required to distinguish field law, representation, and
observation as non-equivalent layers.


\subsection{Angular structure from Maxwell symmetry and boundary conditions}
\label{sec:ch01.0.2}
% TOC tag: [HOME]

Angular organization is approached here through Maxwell boundary-value dynamics,
not through a preliminary list of mode labels \cite{jackson1999,harrington2001}.
Let \(\Omega\subset\mathbb{R}^3\) be the analysis region, \(\Gamma=\partial\Omega\)
its boundary, \(\chi\) a constitutive model descriptor, and \(\Jimp\) the impressed
electric current density. In this setting, the field pair \((\E,\H)\) is viewed
as the image of a constrained Maxwell map
\(\mathcal{M}_{\Omega,\Gamma,\chi}\!\left[\Jimp\right]\). The key consequence is
that angular behavior belongs to operator image structure, not to coordinate
naming alone.

For interpretation, three layers must remain non-interchangeable. Local angular
content appears at density level (\(\rho_J\)), transport through control surfaces
appears at flux level (\(\Phi_J\)), and basis-expansion information appears at
coefficient level (\(\Sigma\)). Phase topology embedded in \(\Sigma\) is therefore
insufficient evidence of transported angular momentum unless \(\Phi_J\) and
conservation accounting are simultaneously satisfied
\cite{jackson1999,harrington2001}.

The symmetry statement governing later chapters follows from the same viewpoint.
When Maxwell action and active rotation commute on the admissible domain,
rotationally adapted organization survives propagation; when commutation is
broken by geometry, boundary anisotropy, or constitutive asymmetry, coupling
between angular channels emerges as a structural necessity
\cite{harrington2001,collin2001}. Chapter~\ref{ch:01} introduces this principle
at interpretation level, while rigorous operator and spectral proofs remain at
their dedicated owner locations in later chapters.


\subsection{Assumptions, finite-energy admissibility, and domain of validity}
\label{sec:ch01.0.3}
% TOC tag: [HOME]

The chapter is developed on a finite-energy admissible class so that inner
products, norm estimates, and projection operations remain both mathematically
well-posed and physically interpretable. For the analysis region
\(\Omega\subset\mathbb{R}^3\), the principal admissibility condition is
\begin{equation}
\label{eq:ch1.energy-class}
\int_{\Omega}\!\left(\E^*\!\cdot\E+\H^*\!\cdot\H\right)\,dV<\infty,
\end{equation}
where \((\cdot)^*\) is complex conjugation and \(dV\) is the volume measure on
\(\Omega\). The integrand terms \(\E^*\!\cdot\E\) and \(\H^*\!\cdot\H\) capture
electric and magnetic field-intensity contributions to norm content. As a full
condition, Eq.~\eqref{eq:ch1.energy-class} excludes infinite-energy states whose
presence would invalidate both measurement interpretation and operator estimates
\cite{jackson1999,collin2001}.

This class is used together with three global assumptions throughout
Chapter~\ref{ch:01}: SI normalization, phasor convention \(\exp(-\jj\omega t)\),
and linear local constitutive closure unless an explicit exception appears in a
later owner location. These assumptions form the compatibility frame that keeps
field law, boundary admissibility, and conservation reasoning inside one model
family.

The validity boundary is explicit. Strong nonlinearity, active gain, rapid
constitutive time modulation, and quantum-field corrections lie outside the
scope of this chapter. They may be referenced only as forward pointers and are
treated formally at designated later chapters
\cite{harrington2001,jackson1999}.

\paragraph{Section 1.0 summary and bridge.}
Section~\ref{sec:ch01.0} has prepared the chapter at three coupled levels:
mission boundary, structural origin of angular behavior, and admissibility
domain. With that preparation complete, Section~\ref{sec:ch01.1} develops the
Maxwell-conservation backbone that governs all subsequent angular arguments in
Chapter~\ref{ch:01}.


\section{Maxwell Theory and Conservation Laws}
\label{sec:ch01.1}

Section~\ref{sec:ch01.1} establishes the governing dynamical core that all later
angular claims must inherit: Maxwell equations in the adopted SI phasor
convention, admissible source and boundary data, well-posed evolution, and the
continuity laws that connect local field structure to measurable transfer of
energy, linear momentum, and angular momentum. The subsection sequence is
organized so that conservation statements are read as consequences of the same
field equations, not as independent postulates.


\subsection{Maxwell equations: differential, integral, vacuum constitutive relations}
\label{sec:ch01.1.1}
% TOC tag: [HOME]

The first technical requirement of Section~\ref{sec:ch01.1} is a single
Maxwell-law statement in the adopted phasor convention, written with enough
detail that later subsections can inherit it without re-expression. Consider a
piecewise smooth region \(\Omega\subset\mathbb{R}^3\) with boundary
\(\Gamma=\partial\Omega\), harmonic factor \(\exp(-\jj\omega t)\) with
\(\omega>0\), phasor fields \(\E=\E(\mathbf{r})\) and \(\H=\H(\mathbf{r})\),
impressed current density \(\Jimp=\Jimp(\mathbf{r})\), and free-charge density
\(\rho=\rho(\mathbf{r})\). Vacuum parameters are \(\varepsilon_0\) and \(\mu_0\)
in SI units \cite{jackson1999,collin2001,harrington2001}.

The local field system is
\begin{equation}
\label{eq:ch1.maxwell-local}
\begin{aligned}
\curl \E&=-\jj\omega\mu_0\H, &
\curl \H&=\Jimp+\jj\omega\varepsilon_0\E,\\
\divg\!\left(\varepsilon_0\E\right)&=\rho, &
\divg\!\left(\mu_0\H\right)&=0.
\end{aligned}
\end{equation}
In Eq.~\eqref{eq:ch1.maxwell-local}, each operand has a distinct role.
The curl operator \(\curl\) extracts local rotational content, while
\(\divg\) measures local flux imbalance. The factor \(-\jj\omega\) converts time
differentiation to harmonic algebra; \(\mu_0\H\) and \(\varepsilon_0\E\) are
flux-scaled magnetic and electric terms; \(\Jimp\) is the externally prescribed
source contribution. Read term by term, the first two relations govern
circulation dynamics, and the second row imposes divergence constraints that make
the system physically consistent.

As a complete system, Eq.~\eqref{eq:ch1.maxwell-local} states that electromagnetic
evolution and source coupling are inseparable from charge and magnetic-flux
constraints. A direct implication appears when \(\divg\) acts on the second curl
relation:
\[
\divg\Jimp-\jj\omega\rho=0.
\]
The first term is source-current divergence and the second is harmonic
charge-rate term; together they enforce phasor charge continuity, preventing
nonconservative source prescriptions.

For oriented surface \(S\subset\Omega\) with boundary \(\partial S\), bounded
volume \(V\subset\Omega\) with boundary \(\partial V\), outward unit normal
\(\hat{\mathbf{n}}\), line element \(d\boldsymbol{\ell}\), surface element \(dS\),
and volume element \(dV\), Stokes and divergence theorems map
Eq.~\eqref{eq:ch1.maxwell-local} to
\begin{equation}
\label{eq:ch1.maxwell-integral}
\begin{aligned}
\oint_{\partial S}\E\cdot d\boldsymbol{\ell}
&=-\jj\omega\mu_0\int_S\H\cdot\hat{\mathbf{n}}\,dS,\\
\oint_{\partial S}\H\cdot d\boldsymbol{\ell}
&=\int_S\Jimp\cdot\hat{\mathbf{n}}\,dS
+\jj\omega\varepsilon_0\int_S\E\cdot\hat{\mathbf{n}}\,dS,\\
\int_{\partial V}\varepsilon_0\E\cdot\hat{\mathbf{n}}\,dS
&=\int_V\rho\,dV,\qquad
\int_{\partial V}\mu_0\H\cdot\hat{\mathbf{n}}\,dS=0.
\end{aligned}
\end{equation}
Operand-level reading of Eq.~\eqref{eq:ch1.maxwell-integral} clarifies
measurement meaning: contour integrals represent circulation response; surface
integrals represent flux through control sets; source and charge terms quantify
injection and storage within those sets. The full equation block gives the finite
control-region form of the same field law, with no new physical postulate added.
Its value is practical and interpretive: quantities appearing in each line map
directly to measurable loops, surfaces, and enclosed volumes.

\begin{figure}[t]
\centering
\fbox{\parbox{0.92\linewidth}{
\textbf{Control-set schematic for Maxwell integral balances.}
An oriented surface \(S\) with boundary \(\partial S\) and an oriented volume \(V\)
with boundary \(\partial V\) are drawn as two complementary measurement windows.
Circulation responses are evaluated on \(\partial S\), flux responses on \(S\) or
\(\partial V\), and free charge is accumulated over \(V\). The geometry makes the
local-to-global transition in Eq.~\eqref{eq:ch1.maxwell-integral}
transparent.}}
\caption{Conceptual geometry linking differential Maxwell laws to circulation and flux measurements in Subsection~\ref{sec:ch01.1.1}.}
\label{fig:ch1.maxwell-control-sets}
\end{figure}

Vacuum closure for flux variables and propagation scale is
\begin{equation}
\label{eq:ch1.vacuum-constitutive}
\mathbf{D}=\varepsilon_0\E,\qquad
\mathbf{B}=\mu_0\H,\qquad
c=\frac{1}{\sqrt{\varepsilon_0\mu_0}},
\end{equation}
where \(\mathbf{D}\) is electric flux density, \(\mathbf{B}\) is magnetic flux
density, and \(c\) is vacuum wave speed. At operand level, \(\varepsilon_0\E\)
and \(\mu_0\H\) convert field intensities to flux densities, while
\(1/\sqrt{\varepsilon_0\mu_0}\) sets intrinsic propagation scale.
As a complete block, Eq.~\eqref{eq:ch1.vacuum-constitutive} closes the vacuum
model in \((\E,\H)\) without extra constitutive operators.

A familiar source-free implication (\(\Jimp=0,\rho=0\)) then appears as
\[
\nabla^2\E+\omega^2\mu_0\varepsilon_0\E=0,\qquad
\nabla^2\H+\omega^2\mu_0\varepsilon_0\H=0,
\]
with \(\nabla^2\) the vector Laplacian. The operator \(\nabla^2\) controls spatial
curvature, and \(\omega^2\mu_0\varepsilon_0\) supplies squared propagation scale.
Physically, free-space fields propagate as constrained wave solutions governed by
the same vacuum constants that set flux closure.

Validity limits are explicit. Equivalence between local and integral forms
requires regularity compatible with Stokes/divergence theorems, and
Eq.~\eqref{eq:ch1.vacuum-constitutive} ceases to suffice for dispersive,
anisotropic, nonlinear, nonlocal, or strongly time-varying media where
constitutive response is operator-valued \cite{jackson1999,collin2001}.

Subsection~\ref{sec:ch01.1.2} can now address source admissibility and boundary
data on this inherited Maxwell base, with no re-expression of the governing
equation set.


\subsection{Sources, boundary conditions, and physical admissibility}
\label{sec:ch01.1.2}
% TOC tag: [HOME]

The Maxwell system of Subsection~\ref{sec:ch01.1.1} becomes predictive only after
its data class is specified with sufficient analytical precision. In this chapter,
\emph{source} means prescribed excitation through impressed electric current
density \(\Jimp\) and free-charge density \(\rho\); \emph{boundary condition}
means a constraint on field traces at \(\Gamma=\partial\Omega\) and at interior
interfaces; \emph{physical admissibility} means joint regularity and compatibility
conditions under which those data produce realizable fields rather than algebraic
extensions without physical content \cite{jackson1999,harrington2001,collin2001}.

The logical role of this subsection in Chapter~\ref{ch:01} is therefore to close
the gap between field equations and usable problem statements. Maxwell equations
provide evolution and constraint structure; admissible data classes decide whether
that structure is well posed for physical interpretation. For readers transitioning
from undergraduate training, the key point is that solution quality depends not
only on the differential equations, but on whether source and boundary data belong
to a mathematically compatible class. In this monograph, that compatibility is
kept explicit because later claims on momentum transport, gauge-sensitive
quantities, and observable/inferred separation all inherit it directly.

Let \(\Omega\subset\mathbb{R}^3\) be decomposed into piecewise smooth material
subdomains \(\{\Omega_a\}_{a=1}^{N_\Omega}\), and let
\(\Sigma_{ab}=\partial\Omega_a\cap\partial\Omega_b\) be an interface with unit
normal \(\hat{\mathbf{n}}_{ab}\) directed from \(\Omega_a\) to \(\Omega_b\). Let
\(\mathbf{D}_a\) and \(\mathbf{B}_a\) denote electric and magnetic flux-density
phasors in \(\Omega_a\), \(\mathbf{J}_{s,ab}\) the electric surface-current
density on \(\Sigma_{ab}\), and \(\rho_{s,ab}\) the free surface-charge density
on \(\Sigma_{ab}\). On \(\Gamma\), data closure is represented by a boundary
operator \(\mathcal{B}_\Gamma\) acting on the boundary traces of \((\E,\H)\):
\begin{equation}
\label{eq:ch1.boundary-data-operator}
\mathcal{B}_\Gamma[\E,\H]=\mathbf{g}_\Gamma,
\end{equation}
where \(\mathbf{g}_\Gamma\) is prescribed boundary data. Equation~\eqref{eq:ch1.boundary-data-operator}
plays two simultaneous roles: mathematically, it closes the boundary-value
problem; physically, it represents how material termination, coating, or idealized
wall behavior restricts admissible tangential traces.

It is useful to write the same closure in domain-range form,
\[
\mathcal{B}_\Gamma:\mathcal{T}_\Gamma\to\mathcal{Y}_\Gamma,
\]
where \(\mathcal{T}_\Gamma\) denotes the trace space generated by admissible
interior fields and \(\mathcal{Y}_\Gamma\) denotes the data space for boundary
specifications. This notation does not introduce a new result; it clarifies that
boundary conditions are operators between function spaces, not merely symbolic
equalities. That viewpoint becomes increasingly important in Chapters~2 and 3,
where operator domain, range, and null-space structure govern accessibility of
angular channels.

Interface transmission follows from the integral Maxwell balances
Eq.~\eqref{eq:ch1.maxwell-integral} by shrinking control loops and pillboxes
across \(\Sigma_{ab}\). For piecewise smooth interfaces and no impressed magnetic
sheet current, one obtains
\begin{equation}
\label{eq:ch1.interface-conditions}
\begin{aligned}
\hat{\mathbf{n}}_{ab}\times(\E_b-\E_a)&=\mathbf{0},\\
\hat{\mathbf{n}}_{ab}\times(\H_b-\H_a)&=\mathbf{J}_{s,ab},\\
\hat{\mathbf{n}}_{ab}\cdot(\mathbf{D}_b-\mathbf{D}_a)&=\rho_{s,ab},\\
\hat{\mathbf{n}}_{ab}\cdot(\mathbf{B}_b-\mathbf{B}_a)&=0.
\end{aligned}
\end{equation}
The analytical structure is explicit. The cross product with \(\hat{\mathbf{n}}_{ab}\)
extracts tangential jumps; the dot product extracts normal-flux jumps. The second
row ties tangential-\(\H\) discontinuity to \(\mathbf{J}_{s,ab}\), and the third
row ties normal-\(\mathbf{D}\) discontinuity to \(\rho_{s,ab}\). Read as a block,
Eq.~\eqref{eq:ch1.interface-conditions} is the transmission law that keeps the
field pair consistent with Maxwell constraints while allowing physically necessary
jump mechanisms.

Intermediate limiting relations leading to Eq.~\eqref{eq:ch1.interface-conditions}
can be read from control-set shrinkage in Eq.~\eqref{eq:ch1.maxwell-integral}:
for a vanishing loop straddling \(\Sigma_{ab}\),
\[
\lim_{\delta\to 0}\oint_{\partial S_\delta}\E\cdot d\boldsymbol{\ell}=0
\quad\Rightarrow\quad
\hat{\mathbf{n}}_{ab}\times(\E_b-\E_a)=\mathbf{0},
\]
and for a vanishing pillbox,
\[
\lim_{\delta\to 0}\int_{\partial V_\delta}\mathbf{D}\cdot\hat{\mathbf{n}}\,dS
=\rho_{s,ab}
\quad\Rightarrow\quad
\hat{\mathbf{n}}_{ab}\cdot(\mathbf{D}_b-\mathbf{D}_a)=\rho_{s,ab}.
\]
These intermediate expressions are included to make the derivational path visible
before the compact final jump form is used in later sections.

From a physical viewpoint, this interface law is the first point where geometry
and material architecture enter as active participants in angular structure:
continuity classes and jump channels determine which field components can carry
transport information across partition lines, and which components are trapped,
screened, or redirected.

Outer closure is specified by partitioning
\(\Gamma=\Gamma_{\mathrm{PEC}}\cup\Gamma_{\mathrm{PMC}}\cup\Gamma_Z\), where
\(\Gamma_{\mathrm{PEC}}\) is perfect-electric-conductor boundary,
\(\Gamma_{\mathrm{PMC}}\) is perfect-magnetic-conductor boundary, and \(\Gamma_Z\)
is an impedance boundary with surface impedance \(Z_s\):
\begin{equation}
\label{eq:ch1.outer-boundary-classes}
\begin{aligned}
\hat{\mathbf{n}}\times\E&=\mathbf{0}\quad &&\text{on }\Gamma_{\mathrm{PEC}},\\
\hat{\mathbf{n}}\times\H&=\mathbf{0}\quad &&\text{on }\Gamma_{\mathrm{PMC}},\\
\hat{\mathbf{n}}\times\E&=Z_s\,(\hat{\mathbf{n}}\times\H)\times\hat{\mathbf{n}}
\quad &&\text{on }\Gamma_Z,
\end{aligned}
\end{equation}
with \(\hat{\mathbf{n}}\) the outward unit normal on \(\Gamma\). Here
\(\hat{\mathbf{n}}\times\E\) and \(\hat{\mathbf{n}}\times\H\) are tangential
electric and magnetic traces, while \((\hat{\mathbf{n}}\times\H)\times\hat{\mathbf{n}}\)
reorients tangential magnetic content into the impedance-coupled direction.
Equation~\eqref{eq:ch1.outer-boundary-classes} therefore functions as a boundary
trace selector mathematically, and as a model of ideal reflection or finite
surface dissipation physically.

For continuity with undergraduate intuition, the three boundary classes may be
viewed as limiting material responses on tangential fields. The PEC relation
suppresses tangential electric trace, the PMC relation suppresses tangential
magnetic trace, and the impedance relation preserves both while coupling them
through \(Z_s\). This hierarchy is not merely classificatory: it determines the
energy-exchange channel at the boundary and therefore influences which angular
content may persist, decay, or be reflected in bounded geometries.

Source admissibility is now imposed by regularity, localization, and continuity
compatibility. Let \(\Omega_s\subseteq\Omega\) be source support. The chapter
source class is
\begin{equation}
\label{eq:ch1.source-admissibility}
\Jimp\in L^2(\Omega_s)^3,\qquad
\operatorname{supp}(\Jimp)\subseteq\overline{\Omega_s},\qquad
\divg\Jimp-\jj\omega\rho=0
\ \text{in }\mathcal{D}'(\Omega).
\end{equation}
The first term in Eq.~\eqref{eq:ch1.source-admissibility} enforces finite source
energy in the local \(L^2\)-sense; the support condition enforces spatial
localization of excitation; the distributional continuity relation enforces
charge-current consistency under the adopted \(\exp(-\jj\omega t)\) phasor
convention. Taken together, these conditions exclude non-conservative source
prescriptions that would violate Eq.~\eqref{eq:ch1.maxwell-local}.

The admissibility class also fixes an implicit domain of validity. If \(\Jimp\) is
too singular for the assumed trace/volume framework, or if charge continuity is
violated by prescribed data, then formal fields can be written but cannot be
interpreted within the same conservation framework used in this chapter. For that
reason, admissibility here is treated as part of physical modeling, not as a
post-processing check.

\begin{proposition}[Global compatibility of source and boundary data]
\label{prop:ch1.source-boundary-compatibility}
Assume Eq.~\eqref{eq:ch1.source-admissibility} and the boundary operator
constraint Eq.~\eqref{eq:ch1.boundary-data-operator}. Then admissible data must
satisfy the global charge-balance identity
\begin{equation}
\label{eq:ch1.global-charge-compatibility}
\int_{\Omega}\!\left(\divg\Jimp-\jj\omega\rho\right)\,dV
=\int_{\Gamma}\!\Jimp\cdot\hat{\mathbf{n}}\,dS-\jj\omega\int_{\Omega}\rho\,dV=0.
\end{equation}
\end{proposition}

The proof is obtained by integrating the distributional continuity relation in
Eq.~\eqref{eq:ch1.source-admissibility} over \(\Omega\), invoking divergence
theorem on \(\divg\Jimp\), and using admissible boundary traces from
Eq.~\eqref{eq:ch1.boundary-data-operator}. The integral
\(\int_{\Gamma}\Jimp\cdot\hat{\mathbf{n}}\,dS\) is net outward impressed-current
flux, while \(\jj\omega\int_{\Omega}\rho\,dV\) is harmonic charge-storage rate.
Equation~\eqref{eq:ch1.global-charge-compatibility} is therefore a global
charge-balance law on admissible data classes.

In practical terms, Eq.~\eqref{eq:ch1.global-charge-compatibility} is the
large-scale consistency test that links boundary current flux to interior charge
storage under the fixed phasor convention. Local constitutive detail may vary
across subdomains, but this global balance must still hold when data are
admissible. That invariance across local complexity is exactly why this equation
is carried as a major milestone relation in the subsection.

\begin{figure}[t]
\centering
\fbox{\parbox{0.94\linewidth}{
\textbf{Interface-boundary-source admissibility map.}
The domain \(\Omega\) is shown with subregions \(\Omega_a,\Omega_b\), interface
\(\Sigma_{ab}\), and exterior boundary partition
\(\Gamma_{\mathrm{PEC}}\cup\Gamma_{\mathrm{PMC}}\cup\Gamma_Z\). Source support
\(\Omega_s\) occupies a compact subset. Interface jumps are constrained by
Eq.~\eqref{eq:ch1.interface-conditions}, boundary traces by
Eq.~\eqref{eq:ch1.outer-boundary-classes}, source regularity by
Eq.~\eqref{eq:ch1.source-admissibility}, and global consistency by
Eq.~\eqref{eq:ch1.global-charge-compatibility}.}}
\caption{Schematic coupling of source support, interface transmission, and outer-boundary closure for admissible Maxwell data in Subsection~\ref{sec:ch01.1.2}.}
\label{fig:ch1.source-boundary-admissibility}
\end{figure}

Figure~\ref{fig:ch1.source-boundary-admissibility} highlights the coupled nature
of admissibility: source support, interface transmission, and exterior trace class
form a single data manifold. Later transport and observable claims inherit this
manifold directly, because angular descriptors computed from \((\E,\H)\) are
credible only when the entire data chain is compatible.

The figure is intentionally schematic rather than geometric-detailed. Its function
is to expose dependency order: source support \(\Omega_s\) fixes where excitation
enters, interface conditions fix how fields pass between material sectors, and
boundary class fixes what traces are allowed to leave or return at \(\Gamma\).
When these three layers are read together, the reader can follow why the chapter
does not treat angular observables before establishing admissibility.

The validity boundary is equally clear. The derivation assumes piecewise smooth
interfaces, local constitutive response within each subregion, and standard trace
theory for boundary operators. Strongly nonlocal media, rough multiscale
interfaces, active time-modulated surfaces, and stochastic boundary operators lie
outside this subsection and require expanded functional frameworks
\cite{jackson1999,harrington2001,collin2001,balanis2016}.

This explicit boundary statement is included for reader confidence: it marks the
exact regime in which the derivation should be trusted and prevents overextension
of conclusions to data classes not yet constructed in the monograph.

This subsection establishes the HOME admissibility frame for Chapter~1 through
Eq.~\eqref{eq:ch1.interface-conditions},
Eq.~\eqref{eq:ch1.outer-boundary-classes},
Eq.~\eqref{eq:ch1.source-admissibility}, and
Proposition~\ref{prop:ch1.source-boundary-compatibility}. With source and trace
classes fixed at this level, Subsection~\ref{sec:ch01.1.3} can address
initial-boundary uniqueness without reopening data admissibility.


\subsection{Initial-boundary value problems and uniqueness}
\label{sec:ch01.1.3}
% TOC tag: [HOME]

With governing equations and admissibility classes already in place, the next
structural task is to specify when a boundary-value prescription yields a single
electromagnetic field pair rather than a family of mathematically admissible but
physically incompatible solutions. The problem data consist of source terms
\((\Jimp,\rho)\) on \(\Omega\), constitutive model \(\chi\), and boundary
operators on \(\Gamma\) chosen from the classes introduced in
Subsection~\ref{sec:ch01.1.2}. For harmonic analysis, uniqueness can be cast in
operator language by asking whether the homogeneous problem admits only the zero
solution.

Let \((\E_1,\H_1)\) and \((\E_2,\H_2)\) be two solutions for the same data and
define the difference fields \(\Delta\E=\E_1-\E_2\),
\(\Delta\H=\H_1-\H_2\). Then \((\Delta\E,\Delta\H)\) satisfies the homogeneous
Maxwell system with homogeneous boundary data. A compact uniqueness criterion is
therefore
\begin{equation}
\label{eq:ch1.uniqueness-null}
\mathcal{B}_{\chi,\Gamma}[\Delta\E,\Delta\H]=0
\Longrightarrow
\Delta\E=\mathbf{0},\ \Delta\H=\mathbf{0},
\end{equation}
where \(\mathcal{B}_{\chi,\Gamma}\) is the boundary-conditioned Maxwell operator
induced by constitutive data \(\chi\) and boundary class \(\Gamma\).
Operand-level reading of Eq.~\eqref{eq:ch1.uniqueness-null} is straightforward:
the left clause expresses vanishing residual of the homogeneous operator with
homogeneous boundary data, and the right clause requires nullity of the field
difference. As a complete statement, Eq.~\eqref{eq:ch1.uniqueness-null} is the
no-ambiguity condition for electromagnetic prediction under fixed admissibility
specification.

The functional mechanism behind Eq.~\eqref{eq:ch1.uniqueness-null} is expressed
through an energy-like coercive estimate on admissible fields:
\begin{equation}
\label{eq:ch1.uniqueness-estimate}
\Re\!\left\langle \mathcal{B}_{\chi,\Gamma}[\Delta\E,\Delta\H],(\Delta\E,\Delta\H)\right\rangle
\ge c_{\chi,\Gamma}\!\left(
\|\Delta\E\|_{L^2(\Omega)^3}^{2}+\|\Delta\H\|_{L^2(\Omega)^3}^{2}
\right),
\end{equation}
with \(c_{\chi,\Gamma}>0\) the coercivity constant tied to material passivity
and boundary dissipation/admissibility class, \(\langle\cdot,\cdot\rangle\) the
problem-compatible sesquilinear pairing, and \(\Re(\cdot)\) the real part.
Each operand has a precise role: the left side measures operator work on the
difference state, while the right side controls the full field-energy norm of
that difference. The full equation expresses that no nonzero homogeneous mode can
persist without violating coercive positivity.

The physical interpretation is central to the chapter logic. Uniqueness is not
an abstract technical ornament; it is the statement that the model does not admit
two different electromagnetic futures under identical source and boundary
prescription. If this condition fails, angular-momentum transport claims become
non-identifiable because distinct fields can satisfy the same external data. When
it holds, later observables and invariants acquire predictive meaning.

\begin{figure}[t]
\centering
\fbox{\parbox{0.92\linewidth}{
\textbf{Conceptual uniqueness map for boundary-value data.}
The same source data \((\Jimp,\rho)\) enter domain \(\Omega\), while boundary
class \(\Gamma\) and constitutive model \(\chi\) define operator
\(\mathcal{B}_{\chi,\Gamma}\). Uniqueness corresponds to collapse of all
homogeneous residual branches into the null branch only, ensuring a single
admissible field response.}}
\caption{Schematic interpretation of uniqueness in Subsection~\ref{sec:ch01.1.3}: fixed data produce one electromagnetic state when the homogeneous null space is trivial.}
\label{fig:ch1.uniqueness-map}
\end{figure}

Validity limits for this subsection are explicit. The estimate
Eq.~\eqref{eq:ch1.uniqueness-estimate} relies on passive constitutive behavior,
boundary classes compatible with energy control, and function-space regularity
consistent with the admissibility framework of Subsection~\ref{sec:ch01.1.2}. In
active media, nonlocal gain-assisted interfaces, or ill-posed inverse boundary
closures, coercivity can fail and uniqueness may be lost
\cite{collin2001,jackson1999,harrington2001}.

With uniqueness architecture established, Subsection~\ref{sec:ch01.1.4} can
develop energy and momentum transport tensors on a solution class whose physical
interpretation is no longer ambiguous.


\subsection{Energy, linear momentum, and Maxwell stress tensor}
\label{sec:ch01.1.4}
% TOC tag: [HOME]

Once existence and uniqueness are secured, transport statements become
interpretable as physical balances rather than formal identities.
Subsection~\ref{sec:ch01.1.4} develops the energy and linear-momentum accounting
framework that later angular-momentum arguments inherit by reference.

In phasor form with the \(\exp(-\jj\omega t)\) convention, let the complex
Poynting vector be \(\mathbf{S}=\tfrac{1}{2}\E\times\H^*\), where
\((\cdot)^*\) is complex conjugation. Let \(u_E=\tfrac{1}{4}\varepsilon_0\E\cdot\E^*\)
and \(u_H=\tfrac{1}{4}\mu_0\H\cdot\H^*\) be electric and magnetic stored-energy
densities, and let \(\Re(\cdot)\), \(\Im(\cdot)\) extract real and imaginary
parts. The complex Poynting theorem may be written as
\begin{equation}
\label{eq:ch1.poynting-complex}
\divg\mathbf{S}
=-\frac{1}{2}\E\cdot\Jimp^*
-\jj\omega\left(u_H-u_E\right).
\end{equation}
Operand-level reading of Eq.~\eqref{eq:ch1.poynting-complex} is decisive.
The divergence operator \(\divg\) on \(\mathbf{S}\) gives net electromagnetic
power outflow density, \(\E\cdot\Jimp^*/2\) is source-field work-rate density,
and \(-\jj\omega(u_H-u_E)\) is reactive exchange between magnetic and electric
storage channels. As a complete equation,
Eq.~\eqref{eq:ch1.poynting-complex} couples power transport, source work, and
local storage imbalance in one conservation-governed statement.

Taking real and imaginary parts yields the standard interpretation split:
\(\Re(\mathbf{S})\) controls net transported power density and
\(\Im(\mathbf{S})\) controls reactive circulation associated with local
near-field energy exchange. This split is essential for the chapter’s angular
program: strong phase structure in reactive regions does not automatically imply
far-field momentum carriage.

Linear momentum flux is represented through the time-averaged Maxwell stress
tensor. In vacuum its dyadic form is
\begin{equation}
\label{eq:ch1.maxwell-stress}
\mathbf{T}
=\frac{1}{2}\Re\!\left[
\varepsilon_0\,\E\E^*
+\mu_0\,\H\H^*
-\frac{1}{2}\mathbf{I}
\left(\varepsilon_0\,\E\cdot\E^*+\mu_0\,\H\cdot\H^*\right)
\right],
\end{equation}
where \(\mathbf{I}\) is the identity dyad and juxtaposition such as \(\E\E^*\)
is dyadic product. Operand by operand, the first two dyads are directional field
intensity contributions, while the isotropic subtraction term
\(\tfrac{1}{2}\mathbf{I}(\cdots)\) removes trace pressure so that directional
tractions are represented correctly. The real-part operator keeps physically
observable time-averaged transfer.

As a complete object, Eq.~\eqref{eq:ch1.maxwell-stress} maps electromagnetic
field state to mechanical traction on material control surfaces. For unit normal
\(\hat{\mathbf{n}}\), the traction density is
\(\mathbf{t}=\mathbf{T}\cdot\hat{\mathbf{n}}\), and the net force on closed
surface \(\partial V\) is \(\mathbf{F}=\int_{\partial V}\mathbf{T}\cdot\hat{\mathbf{n}}\,dS\).
No new numbered equation is introduced for these definitions because they are
direct functional uses of the already-established stress tensor.

\begin{figure}[t]
\centering
\fbox{\parbox{0.92\linewidth}{
\textbf{Energy-momentum control-volume schematic.}
A volume \(V\) with boundary \(\partial V\) receives source work
\(\E\cdot\Jimp^*/2\), exports power via \(\mathbf{S}\), and transmits linear
momentum through boundary traction \(\mathbf{T}\cdot\hat{\mathbf{n}}\). The same
control set therefore supports power-balance interpretation
(Eq.~\eqref{eq:ch1.poynting-complex}) and momentum-flux interpretation
(Eq.~\eqref{eq:ch1.maxwell-stress}).}}
\caption{Conceptual link between power transport and linear-momentum flux in Subsection~\ref{sec:ch01.1.4}.}
\label{fig:ch1.energy-momentum-control}
\end{figure}

The physical implication for the chapter is immediate. Angular interpretations in
later subsections become trustworthy only when built on this energy-momentum
backbone: reactive storage, transported power, and traction-mediated momentum
transfer must be distinguished before any claim about directional or rotational
field organization is accepted.

Validity limits are explicit. Equations~\eqref{eq:ch1.poynting-complex} and
\eqref{eq:ch1.maxwell-stress} are written for vacuum linear local closure and
time-harmonic averaging. In dispersive, anisotropic, or strongly time-varying
media, additional constitutive terms and averaging care are required before using
the same balance form without modification \cite{jackson1999,collin2001}.

Subsection~\ref{sec:ch01.1.5} can now extend this transport framework from linear
momentum to angular momentum and torque balance.


\subsection{Angular momentum conservation and torque balance}
\label{sec:ch01.1.5}
% TOC tag: [HOME]

Subsection~\ref{sec:ch01.1.5} extends the energy-momentum framework of
Subsection~\ref{sec:ch01.1.4} to rotational transport and mechanical torque. The
goal is to identify which field quantities carry angular momentum across control
surfaces and how boundary/material interaction converts that transport into
observable torque.

Let \(\mathbf{r}\) be position vector relative to a declared origin, let
\(\mathbf{T}\) be the Maxwell stress tensor from
Eq.~\eqref{eq:ch1.maxwell-stress}, and let \(\mathbf{t}=\mathbf{T}\cdot\hat{\mathbf{n}}\)
be traction density on boundary \(\partial V\) of control volume \(V\). The
electromagnetic torque density about the declared origin is
\(\boldsymbol{\tau}=\mathbf{r}\times\mathbf{f}\), where \(\mathbf{f}\) is local
force density, and the angular-momentum flux density tensor about that origin is
\(\mathbf{M}=\mathbf{r}\times\mathbf{T}\).

Under these declarations, the local angular-momentum balance is
\begin{equation}
\label{eq:ch1.angular-momentum-balance}
\frac{\partial\mathbf{j}_{\mathrm{em}}}{\partial t}
+\divg\mathbf{M}
=-\boldsymbol{\tau}_{\mathrm{mech}},
\end{equation}
where \(\mathbf{j}_{\mathrm{em}}\) is electromagnetic angular-momentum density
and \(\boldsymbol{\tau}_{\mathrm{mech}}\) is mechanical torque density transferred
to matter. Operand-level roles are explicit: the time-derivative term measures
local storage-rate of electromagnetic angular momentum, the divergence term
\(\divg\mathbf{M}\) measures net angular-momentum outflow from the control set,
and the right side is the sink/source channel that couples field transport to
material torque.

As a complete conservation law, Eq.~\eqref{eq:ch1.angular-momentum-balance}
states that any reduction of field angular momentum in a region must appear as
boundary flux export or mechanical torque deposition. This closes the conceptual
gap between field-structure interpretation and measurable rotational action.
Without this balance relation, phase-structured fields could be misread as
carrying transport-relevant angular momentum even when no torque pathway exists.

In integral form over \(V\), with \(dS\) the surface element on \(\partial V\),
the same law implies
\[
\frac{d}{dt}\int_V \mathbf{j}_{\mathrm{em}}\,dV
+\int_{\partial V}\mathbf{M}\cdot\hat{\mathbf{n}}\,dS
=-\int_V \boldsymbol{\tau}_{\mathrm{mech}}\,dV.
\]
This unnumbered relation is used here as direct functional reading of
Eq.~\eqref{eq:ch1.angular-momentum-balance}, not as a new independent theorem.
The first integral is stored angular-momentum content in the control volume, the
surface integral is transported content across the boundary, and the right-hand
volume integral is net mechanical torque uptake.

\begin{figure}[t]
\centering
\fbox{\parbox{0.92\linewidth}{
\textbf{Torque-balance control-volume schematic.}
A control volume \(V\) around an origin receives and emits angular-momentum flux
\(\mathbf{M}\cdot\hat{\mathbf{n}}\) across \(\partial V\), while material torque
\(\boldsymbol{\tau}_{\mathrm{mech}}\) accumulates inside \(V\). The bookkeeping
connects storage, transport, and torque channels in
Eq.~\eqref{eq:ch1.angular-momentum-balance}.}}
\caption{Conceptual map relating angular-momentum transport to measurable torque in Subsection~\ref{sec:ch01.1.5}.}
\label{fig:ch1.torque-balance-control}
\end{figure}

The physical implication for subsequent chapter sections is fundamental. Angular
labels gain transport meaning only when they remain consistent with both flux and
torque accounting. Fields with pronounced phase winding in reactive zones may
display strong local structure yet produce negligible far-zone torque transfer.
Conversely, radiative configurations with balanced symmetry breaking can deliver
nonzero rotational action even when visual modal descriptors appear modest.

Validity limits must be stated with equal clarity. The balance law in
Eq.~\eqref{eq:ch1.angular-momentum-balance} presumes consistent origin
declaration, admissible stress-tensor framework, and constitutive behavior within
the local linear regime used in this chapter. In strongly nonlocal media, active
metasurfaces, or regimes requiring microscopic spin-lattice exchange models,
additional coupling terms may be necessary before torque partition is interpreted
without ambiguity \cite{jackson1999,collin2001,harrington2001}.

\paragraph{Section 1.1 summary and transition.}
Section~\ref{sec:ch01.1} now forms a closed transport backbone. Maxwell field
law, source-boundary admissibility, uniqueness, energy and linear-momentum
accounting, and angular-momentum torque balance are all fixed at foundational depth. On
this completed base, Section~\ref{sec:ch01.2} examines angular structure in
field-theoretic terms without conflating representational appearance with
conservation-consistent transport.


\section{Angular Structure of Electromagnetic Fields}
\label{sec:ch01.2}

Section~\ref{sec:ch01.2} develops the meaning of angular structure from
field quantities and transport criteria rather than from phase-pattern naming.
Its scope is to separate directional geometry, polarization-phase organization,
and genuine angular-momentum carriage under explicit radiative and reactive
regimes, so that each later modal interpretation remains tied to an operational
electromagnetic criterion.


\subsection{Directional dependence and angular momentum}
\label{sec:ch01.2.1}
% TOC tag: [HOME]

Directional dependence is the first point at which angular description becomes
physically selective. Two fields may carry the same total radiated power while
distributing that power over solid angle in entirely different ways, and only
one of those distributions may support the angular-momentum transport pattern
required by a given interaction geometry. Consequently, angular discourse in this
chapter begins from directional flux content before any representation labels are
introduced \cite{jackson1999,balanis2016}.

Let \(\hat{\mathbf{r}}=\hat{\mathbf{r}}(\theta,\phi)\) denote the radial unit
vector on an observation sphere, and let
\(d\Omega=\sin\theta\,d\theta\,d\phi\) denote solid-angle measure. Using the
complex Poynting vector \(\mathbf{S}\) from Eq.~\eqref{eq:ch1.poynting-complex},
the directional transport density is introduced as
\begin{equation}
\label{eq:ch1.directional-power-density}
\Pi(\theta,\phi;r)=r^2\,\Re\!\left\{\mathbf{S}(r,\theta,\phi)\cdot\hat{\mathbf{r}}\right\},
\end{equation}
where \(\Re\{\cdot\}\) selects time-averaged transport content. The inner product
\(\mathbf{S}\cdot\hat{\mathbf{r}}\) isolates outward transfer, while the factor
\(r^2\) converts local flux density into per-solid-angle flux. In physical terms,
\(\Pi\) is not a decorative angular map; it is the quantity that determines which
directions actually carry electromagnetic action away from the source region.

The scalar total follows by angular integration:
\begin{equation}
\label{eq:ch1.directional-power-integral}
P_{\mathrm{rad}}(r)=\int_{4\pi}\Pi(\theta,\phi;r)\,d\Omega.
\end{equation}
Equation~\eqref{eq:ch1.directional-power-integral} is often treated as the end
of the discussion, but in angular analysis it is only a projection. Distinct
fields may share \(P_{\mathrm{rad}}\) and still differ radically in directional
content, so any claim about angular transport that suppresses
\(\Pi(\theta,\phi;r)\) loses essential structure.

Rotational transport requires the corresponding directional map of
angular-momentum flux. With \(\mathbf{M}\) from
Eq.~\eqref{eq:ch1.angular-momentum-balance}, define
\begin{equation}
\label{eq:ch1.directional-am-flux}
\mathbf{\Psi}_J(\theta,\phi;r)=r^2\,\mathbf{M}(r,\theta,\phi)\cdot\hat{\mathbf{r}}.
\end{equation}
The contraction \(\mathbf{M}\cdot\hat{\mathbf{r}}\) selects outward angular
transport channel-by-channel; \(r^2\) performs the same geometric conversion used
for power flux. The pair
\((\Pi,\mathbf{\Psi}_J)\) provides the minimum directional dataset needed to
disentangle energy carriage from angular-momentum carriage.

For practical interpretation, consider two radiative states with
\(P_{\mathrm{rad}}^{(1)}=P_{\mathrm{rad}}^{(2)}\) but
\(\mathbf{\Psi}_J^{(1)}\neq\mathbf{\Psi}_J^{(2)}\) over part of \(4\pi\). Equal
total power then coexists with unequal rotational transfer potential, so torque
delivery, angular recoil, and receiver-coupling selectivity can differ despite
identical scalar budgets \cite{harrington2001,collin2001}. This resolves a
frequent misconception in which phase-rich appearance is taken as sufficient
evidence of angular transport.

Origin declaration is also inseparable from directional angular interpretation.
Because rotational quantities are moment-based through
\(\mathbf{r}\times(\cdot)\), shifting the origin can redistribute directional
content even when total transport remains fixed. Statements about directional
angular behavior are therefore incomplete unless origin and boundary geometry are
declared with the same precision as source assumptions.

The validity window is restricted to regimes where outward transport
interpretation is meaningful and where surface-flux integrals are not dominated
by near-zone reactive exchange. Reactive-radiative separation will be developed
formally in Subsection~\ref{sec:ch01.2.4}; until then, directional claims are to
be read as propagation-channel statements, not as generic near-field descriptors.

The next subsection enters amplitude, phase, and polarization structure, and does
so on this directional-transport base so that topological description remains
tethered to physically transferable field content.


\subsection{Amplitude, phase, polarization, and Poynting topology}
\label{sec:ch01.2.2}
% TOC tag: [HOME]

Directional transport maps from Subsection~\ref{sec:ch01.2.1} become physically
informative only after resolving how amplitude, phase, and polarization are
embedded in the complex field pair. For a local transverse chart
\((\hat{\boldsymbol{\theta}},\hat{\boldsymbol{\phi}})\), write
\begin{equation}
\label{eq:ch1.transverse-components}
\E_t=E_\theta\,\hat{\boldsymbol{\theta}}+E_\phi\,\hat{\boldsymbol{\phi}},\qquad
E_\theta=A_\theta e^{\jj\varphi_\theta},\quad
E_\phi=A_\phi e^{\jj\varphi_\phi},
\end{equation}
where \(A_\theta,A_\phi\ge 0\) are amplitude envelopes and
\(\varphi_\theta,\varphi_\phi\) are component phases. Equation
\eqref{eq:ch1.transverse-components} separates magnitude and phase channels at
operand level, which is necessary before any statement about polarization or
topological energy flow can be considered transport-relevant.

The physically active phase variable is relative phase,
\(\Delta\varphi=\varphi_\phi-\varphi_\theta\), not absolute phase alone. A
compact polarization invariant set is obtained through Stokes-type quantities:
\begin{equation}
\label{eq:ch1.stokes-local}
\begin{aligned}
S_0&=|E_\theta|^2+|E_\phi|^2,\\
S_1&=|E_\theta|^2-|E_\phi|^2,\\
S_2&=2\Re\!\left\{E_\theta E_\phi^*\right\},\\
S_3&=-2\Im\!\left\{E_\theta E_\phi^*\right\},
\end{aligned}
\end{equation}
with \((\cdot)^*\) denoting complex conjugation and \(\Im\{\cdot\}\) imaginary
part. In Eq.~\eqref{eq:ch1.stokes-local}, \(S_0\) measures local intensity,
\(S_1,S_2\) encode linear-polarization bias, and \(S_3\) encodes handedness
content. Since \(S_2,S_3\) depend on \(E_\theta E_\phi^*\), polarization state
is phase-coupled; amplitude-only descriptions are therefore incomplete for
angular interpretation \cite{jackson1999,balanis2016}.

Transport interpretation enters through Poynting topology. Using
Eq.~\eqref{eq:ch1.poynting-complex}, define time-averaged energy-flow field
\begin{equation}
\label{eq:ch1.poynting-average}
\mathbf{P}=\Re\!\left\{\mathbf{S}\right\},
\end{equation}
and local flow vorticity indicator
\begin{equation}
\label{eq:ch1.poynting-vorticity}
\boldsymbol{\Omega}_P=\curl \mathbf{P}.
\end{equation}
Equation~\eqref{eq:ch1.poynting-average} isolates transported power density,
while Eq.~\eqref{eq:ch1.poynting-vorticity} diagnoses local circulation tendency
of that transport field. Regions with strong \(\boldsymbol{\Omega}_P\) do not
automatically imply net angular-momentum export; they indicate structured local
flow geometry whose transport significance must still be tested against the
angular flux criteria fixed earlier.

Phase singular points, where both transverse components vanish,
\begin{equation}
\label{eq:ch1.phase-singularity-condition}
A_\theta=A_\phi=0,
\end{equation}
form topological loci around which phase winds are defined. Equation
\eqref{eq:ch1.phase-singularity-condition} marks geometry of phase defects, but
topological charge near such loci is not by itself a torque certificate.
Without nonzero outward angular flux and conservation-consistent channel balance,
singularity structure remains a representation feature rather than a guaranteed
transport outcome \cite{harrington2001,collin2001}.

An allied implication is crucial for this monograph’s interpretive discipline.
Two fields may share nearly identical intensity map \(S_0\) and yet carry
different \(\Delta\varphi\), different \(S_3\), and different
\(\boldsymbol{\Omega}_P\); their coupling to rotational observables can then
diverge sharply despite similar scalar energy appearance. Amplitude map,
polarization state, and flow topology must therefore be read as a coupled
triplet when discussing angular structure.

Validity limits remain explicit. The local decomposition above presumes
sufficient field regularity for component extraction and differentiability of
\(\mathbf{P}\) where \(\curl\mathbf{P}\) is interpreted. In highly lossy,
strongly inhomogeneous, or measurement-truncated settings, topology indicators
require controlled windowing and bandwidth qualification before cross-section
comparison.

With amplitude-phase-polarization-topology coupling now fixed, the next
subsection states the necessary conditions under which this structured field
organization results in genuine angular-momentum transport.


\subsection{Conditions for angular-momentum transport}
\label{sec:ch01.2.3}
% TOC tag: [HOME]

A field configuration exhibits genuine angular-momentum transport only when
directional structure, conservation balance, and outward channel support hold
simultaneously. Any one of these in isolation is insufficient. The diagnostic
question is therefore not whether angular patterns are visible, but whether
those patterns survive flux accounting and can be exported through admissible
propagation channels.

Let \(\mathbf{\Psi}_J(\theta,\phi;r)\) be the directional angular-momentum flux
map from Eq.~\eqref{eq:ch1.directional-am-flux}, and let
\(\Pi(\theta,\phi;r)\) be directional power flux from
Eq.~\eqref{eq:ch1.directional-power-density}. A necessary directional condition
for nontrivial net transport over a chosen observation sphere is
\begin{equation}
\label{eq:ch1.am-net-transport-condition}
\int_{4\pi}\mathbf{\Psi}_J(\theta,\phi;r)\,d\Omega \neq \mathbf{0}.
\end{equation}
Equation~\eqref{eq:ch1.am-net-transport-condition} states that local angular
structure must accumulate into a nonzero exported balance; otherwise the pattern
is confined to internal recirculation or compensating directional cancellation.

Directional sufficiency, however, cannot be judged from
Eq.~\eqref{eq:ch1.am-net-transport-condition} alone. The same field must satisfy
global conservation consistency with the local law already fixed in
Eq.~\eqref{eq:ch1.angular-momentum-balance}. For a control volume \(V\) with
boundary \(\partial V\), this consistency is represented by
\begin{equation}
\label{eq:ch1.am-control-balance}
\frac{d}{dt}\int_V \mathbf{j}_{\mathrm{em}}\,dV
+\int_{\partial V}\mathbf{M}\cdot\hat{\mathbf{n}}\,dS
=-\int_V \boldsymbol{\tau}_{\mathrm{mech}}\,dV,
\end{equation}
where \(\mathbf{j}_{\mathrm{em}}\) is electromagnetic angular-momentum density,
\(\mathbf{M}\) is angular-momentum flux tensor, \(\hat{\mathbf{n}}\) is outward
unit normal on \(\partial V\), and \(\boldsymbol{\tau}_{\mathrm{mech}}\) is
mechanical torque density delivered to matter. Equation
\eqref{eq:ch1.am-control-balance} rules out interpretive shortcuts: exported
angular transport must appear in torque or storage channels, not in phase
appearance alone \cite{jackson1999,harrington2001}.

Transport classification also requires separation of outward propagative support
from reactive exchange. Using \(\mathbf{P}=\Re\{\mathbf{S}\}\) from
Eq.~\eqref{eq:ch1.poynting-average}, outward channel support for a directional
sector \(\Delta\Omega\subseteq 4\pi\) is checked by
\begin{equation}
\label{eq:ch1.propagative-support}
\int_{\Delta\Omega}\Pi(\theta,\phi;r)\,d\Omega>0.
\end{equation}
Equation~\eqref{eq:ch1.propagative-support} enforces that angular-momentum
claims are tied to sectors carrying net outward power; in sectors where this
condition fails, strong local field complexity may persist without transport
credibility.

Three criteria therefore define the transport-ready regime in this subsection:
nonzero directional angular export
\eqref{eq:ch1.am-net-transport-condition}, conservation-consistent control
balance \eqref{eq:ch1.am-control-balance}, and outward propagative support
\eqref{eq:ch1.propagative-support}. Field states violating any one criterion
remain representation-rich but transport-incomplete. This triad prevents a
common error in advanced OAM discourse, where topological or modal appearance is
mistaken for torque-capable transport without flux validation
\cite{collin2001,balanis2016}.

Origin declaration remains part of the condition set because angular moments are
computed through \(\mathbf{r}\times(\cdot)\). Changing origin can redistribute
directional partitions and alter modal bookkeeping; transport conclusions are
therefore comparable only when origin and boundary geometry are fixed across the
comparison set.

In practice, these conditions define the admissible interpretation pathway for
the rest of Chapter~\ref{ch:01}: first establish outward channel support, then
verify conservation consistency, and only thereafter interpret directional phase
and polarization structure as angular transport evidence.

The remaining distinction between reactive storage and radiative export is now
taken up explicitly in Subsection~\ref{sec:ch01.2.4}, where regime boundaries
are formalized.


\subsection{Reactive and radiative regimes}
\label{sec:ch01.2.4}
% TOC tag: [HOME]

A physically credible angular-transport claim must distinguish energy that is
exported irreversibly from energy that is stored and cyclically exchanged in
the source neighborhood. The distinction is not semantic: it controls whether
the directional structure identified in
Eq.~\eqref{eq:ch1.directional-am-flux} represents torque-capable propagation or
only near-field complexity that does not survive far-field accounting.

Assume a time-harmonic source in a linear, passive medium region outside a
bounded excitation volume, with phasor convention fixed by
\(e^{-j\omega t}\). Let \(\mathbf{S}(\mathbf{r})\) denote the complex Poynting
vector already introduced in Eq.~\eqref{eq:ch1.poynting-complex}, let
\(\hat{\mathbf{r}}\) denote the radial unit vector on a sphere of radius \(r\),
and let \(d\Omega\) denote the solid-angle element. The net outward active power
through that sphere is
\begin{equation}
\label{eq:ch1.rad-power-sphere}
P_{\mathrm{rad}}(r)=r^2\int_{4\pi}\Re\!\left\{\mathbf{S}\!\cdot\!\hat{\mathbf{r}}\right\}\,d\Omega,
\end{equation}
where \(\Re\{\cdot\}\) extracts real part and therefore retains only irreversible
transport contribution. The companion reactive exchange indicator is defined by
\begin{equation}
\label{eq:ch1.reactive-flux-sphere}
Q_{\mathrm{re}}(r)=r^2\int_{4\pi}\Im\!\left\{\mathbf{S}\!\cdot\!\hat{\mathbf{r}}\right\}\,d\Omega,
\end{equation}
where \(\Im\{\cdot\}\) extracts imaginary part associated with oscillatory
electric-magnetic energy interchange over one period
\cite{harrington2001,collin2001}.

The regime classification used in this chapter follows from how
\(P_{\mathrm{rad}}(r)\) and \(Q_{\mathrm{re}}(r)\) scale with observation
radius. A transport-supporting radiative regime is identified when
\begin{equation}
\label{eq:ch1.radiative-regime-condition}
\lim_{r\to\infty}P_{\mathrm{rad}}(r)=P_\infty>0,\qquad
\lim_{r\to\infty}\frac{|Q_{\mathrm{re}}(r)|}{P_{\mathrm{rad}}(r)}=0,
\end{equation}
with \(P_\infty\) the asymptotic emitted power. The first limit demands
non-vanishing outward export; the second limit suppresses reactive dominance at
large radius. When this pair fails, the field may still display substantial
local phase or polarization organization, but that organization cannot be read
as robust radiative transport.

For finite-radius diagnosis, which is the experimentally relevant setting,
define the dimensionless regime index
\begin{equation}
\label{eq:ch1.regime-index}
\chi(r)=\frac{|Q_{\mathrm{re}}(r)|}{P_{\mathrm{rad}}(r)+\epsilon_P},
\end{equation}
where \(\epsilon_P>0\) is a small regularization constant introduced only to
avoid division singularity in weak-radiation windows. Values \(\chi(r)\gg1\)
indicate reactive dominance, while \(\chi(r)\ll1\) indicates export-dominated
behavior. Equation~\eqref{eq:ch1.regime-index} does not replace the asymptotic
criterion in Eq.~\eqref{eq:ch1.radiative-regime-condition}; it operationalizes
the same separation when measurements are restricted to finite apertures.

This separation immediately refines Subsection~\ref{sec:ch01.2.3}. The
directional transport condition of Eq.~\eqref{eq:ch1.am-net-transport-condition}
and the control-balance requirement of Eq.~\eqref{eq:ch1.am-control-balance}
remain necessary, yet they are interpreted differently depending on \(\chi(r)\):
for reactive-dominant shells, angular patterns mostly diagnose storage geometry;
for radiative-dominant shells, the same patterns become credible indicators of
exportable angular momentum. The mathematical partition therefore resolves a
common ambiguity in advanced field interpretation, namely conflating visible
helical/azimuthal structure with radiative delivery in all regimes
\cite{jackson1999,balanis2016}.

The validity boundary of this subsection is explicit. The scalar index
\(\chi(r)\) and the limits in Eq.~\eqref{eq:ch1.radiative-regime-condition}
presume linear media and time-harmonic steady state. Broadband pulses, strongly
dispersive media with nonlocal response, or strongly nonlinear conversion can
require windowed or generalized energy-partition operators not introduced in
Chapter~\ref{ch:01}. Within the stated assumptions, however, the
reactive-radiative distinction is sufficient to prevent regime-mixing errors in
all later angular-structure arguments.

With regime separation now fixed at foundational level, Section~\ref{sec:ch01.3} can
treat rotational covariance without conflating coordinate effects with near-field
storage effects.

\paragraph{Section 1.2 synthesis and bridge.}
Section~\ref{sec:ch01.2} identifies the conditions under which angular features
correspond to transport-capable electromagnetic content. This clarification
prepares the covariance analysis in Section~\ref{sec:ch01.3}, where rotational
actions and tensor structure are treated explicitly.


\section{Rotations, Tensors, and Covariance}
\label{sec:ch01.3}

Section~\ref{sec:ch01.3} places angular discourse in the language of coordinate
transformations and invariant statements. The objective is to ensure that claims
about angular behavior remain physically interpretable under frame changes,
origin shifts, and tensorial reformulation, thereby preventing coordinate
artifacts from entering theorem-level conclusions.


\subsection{Active and passive rotations}
\label{sec:ch01.3.1}
% TOC tag: [HOME]

Rotational analysis in this chapter pivots on a distinction that is simple to
state yet easy to violate in advanced calculations: rotating the field itself is
not the same operation as rotating only the coordinate chart used to describe
that field. The difference is the first safeguard against attributing
representation changes to physics.

Let \(R\in\mathrm{SO}(3)\) be a proper rotation, with \(R^\top R=I\),
\(\det R=1\), and \(I\) the \(3\times3\) identity matrix. For the phasor electric
field \(\E(\mathbf{r})\) at position \(\mathbf{r}\in\mathbb{R}^3\), active action
is
\begin{equation}
\label{eq:ch1.active-rotation-action}
\big(\mathcal{A}_R\E\big)(\mathbf{r})=R\,\E\!\left(R^{-1}\mathbf{r}\right),
\end{equation}
where \(\mathcal{A}_R\) is the active-rotation operator. The argument
\(R^{-1}\mathbf{r}\) performs geometric pullback to the pre-rotation point, while
left multiplication by \(R\) rotates the vector components themselves. Equation
\eqref{eq:ch1.active-rotation-action} therefore captures rigid pattern
reorientation in physical space \cite{jackson1999,harrington2001}.

The passive counterpart keeps the physical configuration fixed and changes only
the coordinate description:
\begin{equation}
\label{eq:ch1.passive-rotation-action}
\big(\mathcal{P}_R\E\big)(\mathbf{r})=R^{-1}\,\E\!\left(R\mathbf{r}\right),
\end{equation}
with \(\mathcal{P}_R\) the passive-rotation operator. Active and passive
descriptions are connected exactly by
\begin{equation}
\label{eq:ch1.active-passive-duality}
\mathcal{A}_R=\mathcal{P}_{R^{-1}}.
\end{equation}
The algebra in Eq.~\eqref{eq:ch1.active-passive-duality} asserts that either
description is valid if inverse group elements are paired correctly.

\begin{figure}[t]
\centering
\fbox{\begin{minipage}{0.92\linewidth}
\centering
\textbf{Conceptual rotation map for Eq.~\eqref{eq:ch1.active-rotation-action} and Eq.~\eqref{eq:ch1.passive-rotation-action}}\\[2mm]
Node A: physical field pattern before rotation, \(\E\).\\
Node B: actively rotated pattern, \(\mathcal{A}_R\E\).\\
Node C: same physical pattern as A, represented in rotated coordinates, \(\mathcal{P}_R\E\).\\[1mm]
Arrow A\(\rightarrow\)B: rigid geometric action \(R\).\\
Arrow A\(\rightarrow\)C: basis/chart relabeling \(R^{-1}\).\\
Consistency statement: B and C encode the same geometry through inverse group actions.
\end{minipage}}
\caption{Active-passive rotation duality schematic. Labels A, B, and C identify the three representation states used in Eqs.~\eqref{eq:ch1.active-rotation-action}--\eqref{eq:ch1.active-passive-duality}.}
\label{fig:ch1.active-passive-rotation}
\end{figure}

Figure~\ref{fig:ch1.active-passive-rotation} gives the interpretation layer for
the three equations above. The A\(\rightarrow\)B arrow corresponds to genuine
physical reorientation; the A\(\rightarrow\)C arrow corresponds to relabeling of
components at fixed geometry. The duality relation states that these two routes
must agree when the inverse action is used in the passive channel.

Observable balance provides the physical verification. With \(\mathbf{S}\) the
complex Poynting vector from Eq.~\eqref{eq:ch1.poynting-complex},
\(\hat{\mathbf{n}}\) the unit outward normal on sphere \(S_r\) of radius \(r\),
and \(dS\) the surface element, emitted active power remains invariant:
\begin{equation}
\label{eq:ch1.rotation-power-invariance}
\int_{S_r}\Re\!\left\{\mathbf{S}_R\cdot\hat{\mathbf{n}}\right\}\,dS
=
\int_{S_r}\Re\!\left\{\mathbf{S}\cdot\hat{\mathbf{n}}\right\}\,dS.
\end{equation}
Here \(\mathbf{S}_R\) is the transformed complex Poynting field and
\(\Re\{\cdot\}\) extracts active transport contribution. Equation
\eqref{eq:ch1.rotation-power-invariance} prevents a common interpretive error:
directional redistribution under rotation is admissible, but net radiative export
cannot be changed by rigid frame action, in agreement with
Eq.~\eqref{eq:ch1.radiative-regime-condition} \cite{collin2001,balanis2016}.

For generator-level structure, consider a small angle \(\delta\varphi\) about
unit axis \(\hat{\mathbf{u}}\), \(\|\hat{\mathbf{u}}\|=1\):
\begin{equation}
\label{eq:ch1.infinitesimal-active-rotation}
\big(\mathcal{A}_{\delta\varphi}\E\big)(\mathbf{r})
=
\E(\mathbf{r})
+\delta\varphi\!\left[\hat{\mathbf{u}}\times\E(\mathbf{r})
-\big(\hat{\mathbf{u}}\times\mathbf{r}\big)\cdot\nabla\E(\mathbf{r})\right]
+O(\delta\varphi^2),
\end{equation}
where \(\nabla\E\) denotes componentwise spatial derivative and
\(O(\delta\varphi^2)\) collects second-order terms. The first correction rotates
local orientation at fixed point; the second shifts sampling location along the
rotation orbit. Their coexistence explains why local polarization change alone is
insufficient to represent rotational action on fields.

The derivation above is valid for rigid Euclidean rotations in orthonormal frames.
Non-rigid deformations, curvilinear manifolds with connection terms, and
orientation-dependent constitutive manifolds require extensions deferred later.

With active-passive structure fixed at foundational level and interpreted through
Fig.~\ref{fig:ch1.active-passive-rotation}, the next subsection formalizes the
transformation rules for vectors, second-rank tensors, and dyadic maps.


\subsection{Transformation of vectors, tensors, and dyads}
\label{sec:ch01.3.2}
% TOC tag: [HOME]

Covariance becomes computationally stable only after each algebraic object is
assigned its own transformation rule. In angular electromagnetics, failures at
this step propagate directly into false claims about stress transfer, torque
balance, and modal coupling.

Assume the same proper rotation \(R\in\mathrm{SO}(3)\) used in
Subsection~\ref{sec:ch01.3.1}. Let \(\mathbf{v}=v_j\hat{\mathbf{e}}_j\) be a
vector in orthonormal basis \(\hat{\mathbf{e}}_j\), with repeated indices summed
over \(j=1,2,3\). Component transformation is
\begin{equation}
\label{eq:ch1.vector-rotation-components}
v_i' = R_{ij}v_j,
\end{equation}
where \(R_{ij}\) are rotation-matrix elements and \(v_i'\) are rotated
components. Equation~\eqref{eq:ch1.vector-rotation-components} is the finite
dimensional component realization of Eq.~\eqref{eq:ch1.active-rotation-action}.

For a second-rank tensor \(\mathbf{T}\) with components \(T_{pq}\), covariance of
both slots yields
\begin{equation}
\label{eq:ch1.tensor-rotation-rule}
\mathbf{T}' = R\,\mathbf{T}\,R^\top,
\qquad
T'_{ij}=R_{ip}R_{jq}T_{pq}.
\end{equation}
The factor \(R_{ip}\) acts on output index \(p\to i\), while \(R_{jq}\) acts on
input index \(q\to j\). In physical terms, this law guarantees that the Maxwell
stress tensor from Eq.~\eqref{eq:ch1.maxwell-stress} rotates as a geometric flux
object rather than as a fixed numerical table \cite{jackson1999,collin2001}.

Dyadic maps obey the same conjugation pattern. If \(\mathbf{D}\) maps
\(\mathbf{u}\) to \(\mathbf{w}\) through \(\mathbf{w}=\mathbf{D}\cdot\mathbf{u}\),
then
\begin{equation}
\label{eq:ch1.dyadic-covariance}
\mathbf{D}' = R\,\mathbf{D}\,R^\top,
\qquad
\mathbf{w}'=\mathbf{D}'\cdot\mathbf{u}'=R\mathbf{w},
\end{equation}
with \(\mathbf{u}'=R\mathbf{u}\). The first relation transforms the operator, and
the second confirms equivariance of the input-output chain under the same group
action.

\begin{figure}[t]
\centering
\fbox{\begin{minipage}{0.92\linewidth}
\centering
\textbf{Covariance chain for vector/tensor/dyadic transformation}\\[2mm]
Level V (vector): \(\mathbf{v}\xrightarrow{R}\mathbf{v}'\).\\
Level T (tensor): \(\mathbf{T}\xrightarrow{R(\cdot)R^\top}\mathbf{T}'\).\\
Level D (dyadic map): \(\mathbf{u}\xrightarrow{\mathbf{D}}\mathbf{w}\), then
\(\mathbf{u}'=R\mathbf{u}\xrightarrow{\mathbf{D}'}\mathbf{w}'=R\mathbf{w}\).\\[1mm]
Traction branch: \(\mathbf{t}=\mathbf{T}\cdot\hat{\mathbf{n}}\),
\(\hat{\mathbf{n}}'=R\hat{\mathbf{n}}\), \(\mathbf{t}'=\mathbf{T}'\cdot\hat{\mathbf{n}}'\).
\end{minipage}}
\caption{Schematic ladder linking Eqs.~\eqref{eq:ch1.vector-rotation-components}, \eqref{eq:ch1.tensor-rotation-rule}, and \eqref{eq:ch1.dyadic-covariance}. The traction branch indicates how tensor covariance feeds measurable surface force direction.}
\label{fig:ch1.tensor-dyadic-covariance}
\end{figure}

Figure~\ref{fig:ch1.tensor-dyadic-covariance} organizes the three transformation
levels into a single readable path. Level V defines component rotation, Level T
extends the action to two-index objects, and Level D enforces compatibility of a
linear map with rotated input-output vectors. The traction branch is the direct
bridge from covariance algebra to measurable mechanical effect.

Using \(\hat{\mathbf{n}}\) for unit surface normal and
\(\mathbf{t}=\mathbf{T}\cdot\hat{\mathbf{n}}\) for traction, with
\(\hat{\mathbf{n}}'=R\hat{\mathbf{n}}\), one obtains
\begin{equation}
\label{eq:ch1.traction-rotation-covariance}
\mathbf{t}'=\mathbf{T}'\cdot\hat{\mathbf{n}}' = R\mathbf{t}.
\end{equation}
Equation~\eqref{eq:ch1.traction-rotation-covariance} secures the physical meaning
of the figure: transformed traction is exactly the rotated original traction.
Therefore the conservation results from Section~\ref{sec:ch01.1} and angular
transport diagnostics from Section~\ref{sec:ch01.2} remain objective under rigid
coordinate action \cite{harrington2001,balanis2016}.

The derivation is restricted to orthonormal Euclidean frames and proper rotations.
Improper transformations, curvilinear coordinates with connection coefficients,
and material manifolds with intrinsic anisotropic geometry are beyond the present
scope.

With vector, tensor, and dyadic covariance fixed and visually consolidated in
Fig.~\ref{fig:ch1.tensor-dyadic-covariance}, the next subsection isolates origin
dependence from genuine translational covariance in angular measures.


\subsection{Origin dependence and translational covariance}
\label{sec:ch01.3.3}
% TOC tag: [HOME]

Origin choice enters angular diagnostics through the explicit position factor in
moment definitions. If that dependence is not separated from genuine transport,
coordinate bookkeeping can be mistaken for physics. The objective of this
subsection is to derive the exact origin-shift law and then extract a translation
covariant observable core.

Let \(O\) be a reference origin and let \(O'\) be shifted by vector
\(\mathbf{a}\in\mathbb{R}^3\), so that \(\mathbf{r}'=\mathbf{r}-\mathbf{a}\),
where \(\mathbf{r}\) and \(\mathbf{r}'\) are position vectors relative to \(O\)
and \(O'\), respectively. Let \(\mathbf{p}_{\mathrm{em}}(\mathbf{r})\) be
electromagnetic linear-momentum density. Consistent with
Eq.~\eqref{eq:ch1.poynting-average}, one may use
\(\mathbf{p}_{\mathrm{em}}=\mathbf{P}/c^2\), with \(\mathbf{P}\) the active
Poynting vector and \(c\) light speed in the medium approximation used here.
Define total angular momentum in volume \(V\) by
\begin{equation}
\label{eq:ch1.origin-angular-definition}
\mathbf{J}_O=\int_V \mathbf{r}\times\mathbf{p}_{\mathrm{em}}\,dV,
\qquad
\mathbf{J}_{O'}=\int_V \mathbf{r}'\times\mathbf{p}_{\mathrm{em}}\,dV.
\end{equation}
The cross product \(\mathbf{r}\times\mathbf{p}_{\mathrm{em}}\) gives local moment
arm times momentum density; integration accumulates total angular content in the
chosen control region.

Substituting \(\mathbf{r}'=\mathbf{r}-\mathbf{a}\) into
Eq.~\eqref{eq:ch1.origin-angular-definition} yields
\begin{equation}
\label{eq:ch1.origin-shift-law}
\mathbf{J}_{O'}=\mathbf{J}_O-\mathbf{a}\times\mathbf{P}_{\mathrm{tot}},
\qquad
\mathbf{P}_{\mathrm{tot}}=\int_V \mathbf{p}_{\mathrm{em}}\,dV.
\end{equation}
Equation~\eqref{eq:ch1.origin-shift-law} is the central result: origin shift
adds an extrinsic term determined by total linear momentum
\(\mathbf{P}_{\mathrm{tot}}\). Mathematically, this is affine covariance of
angular moments; physically, it means identical fields may report different
\(\mathbf{J}\) values when origins differ, even though transport content is
unchanged.

\begin{figure}[t]
\centering
\fbox{\begin{minipage}{0.92\linewidth}
\centering
\textbf{Origin-shift decomposition of angular momentum}\\[2mm]
Frame A: origin \(O\), moment \(\mathbf{J}_O\).\\
Frame B: origin \(O'=O+\mathbf{a}\), moment \(\mathbf{J}_{O'}\).\\
Link term: \(-\mathbf{a}\times\mathbf{P}_{\mathrm{tot}}\) from
Eq.~\eqref{eq:ch1.origin-shift-law}.\\
Invariant core: intrinsic part after removal of extrinsic translation term.
\end{minipage}}
\caption{Conceptual map of origin dependence: total angular momentum changes by an extrinsic translation term while intrinsic content remains covariant.}
\label{fig:ch1.origin-shift-map}
\end{figure}

Figure~\ref{fig:ch1.origin-shift-map} clarifies the interpretation of
Eq.~\eqref{eq:ch1.origin-shift-law}: moving the origin does not alter field
dynamics; it alters the lever arm. Consequently, comparison across experiments or
simulations is meaningful only when the same origin protocol is enforced or when
extrinsic translation contribution is removed.

A translation-covariant intrinsic observable is obtained by referencing a declared
centroid \(\mathbf{r}_c\) of the measurement protocol:
\begin{equation}
\label{eq:ch1.intrinsic-angular-observable}
\mathbf{J}_{\mathrm{int}}
=
\mathbf{J}_O-\mathbf{r}_c\times\mathbf{P}_{\mathrm{tot}}.
\end{equation}
If both origin and centroid are shifted by \(\mathbf{a}\), then
\(\mathbf{J}_{\mathrm{int}}\) remains unchanged by direct substitution of
Eq.~\eqref{eq:ch1.origin-shift-law}. Operationally,
Eq.~\eqref{eq:ch1.intrinsic-angular-observable} gives a report format compatible
with channel definitions in later sections and consistent with the objectivity
tests developed in Section~\ref{sec:ch01.3}.

The derivation assumes finite momentum integral, regular control boundaries, and
rigid Euclidean translations. Open systems with unresolved momentum flux at
infinity, or media with nonlocal constitutive coupling, require extended flux
closure before Eq.~\eqref{eq:ch1.origin-shift-law} can be applied without
qualification.

With origin dependence isolated and intrinsic observable form fixed, the final
subsection of Section~\ref{sec:ch01.3} can state frame-independent observability
criteria without conflating translation artifacts with physical angular content.


\subsection{Invariants and frame-independent observables}
\label{sec:ch01.3.4}
% TOC tag: [HOME]

Physical interpretation remains stable across coordinate changes only when the
reported quantity is either invariant as a scalar or covariant in a way that
preserves its geometric class. The analysis below establishes an operational test
for that distinction and ties it to the angular-transport diagnostics already
introduced in Section~\ref{sec:ch01.2}.

Let \(\mathcal{F}\) denote the full electromagnetic state in the active-domain
sense of Eq.~\eqref{eq:ch1.active-rotation-action}, and let
\(\mathcal{A}_R\mathcal{F}\) be its rotated counterpart under
\(R\in\mathrm{SO}(3)\). A scalar observable functional
\(\mathcal{I}[\mathcal{F}]\in\mathbb{R}\) is frame-independent if
\begin{equation}
\label{eq:ch1.scalar-observable-invariance}
\mathcal{I}[\mathcal{A}_R\mathcal{F}] = \mathcal{I}[\mathcal{F}]
\quad \forall R\in\mathrm{SO}(3),
\end{equation}
where \(\forall\) means for every admissible proper rotation. A vector-valued
observable \(\mathbf{Q}[\mathcal{F}]\in\mathbb{R}^3\) is frame-objective when
\begin{equation}
\label{eq:ch1.vector-observable-covariance}
\mathbf{Q}[\mathcal{A}_R\mathcal{F}] = R\,\mathbf{Q}[\mathcal{F}].
\end{equation}
Equation~\eqref{eq:ch1.scalar-observable-invariance} characterizes invariance;
Eq.~\eqref{eq:ch1.vector-observable-covariance} characterizes covariance without
loss of geometric identity. Any reported quantity violating both relations is a
representation artifact rather than a physical observable \cite{jackson1999}.

Two scalar invariants extracted from a second-rank tensor
\(\mathbf{T}[\mathcal{F}]\), transformed as in
Eq.~\eqref{eq:ch1.tensor-rotation-rule}, are
\begin{equation}
\label{eq:ch1.tensor-principal-invariants}
I_1(\mathbf{T})=\mathrm{tr}(\mathbf{T}),
\qquad
I_2(\mathbf{T})=\frac{1}{2}\!\left[(\mathrm{tr}\,\mathbf{T})^2-\mathrm{tr}(\mathbf{T}^2)\right],
\end{equation}
where \(\mathrm{tr}(\cdot)\) denotes trace. Because trace is invariant under
similarity action \(R\mathbf{T}R^\top\), both \(I_1\) and \(I_2\) satisfy
Eq.~\eqref{eq:ch1.scalar-observable-invariance}. For the Maxwell stress tensor
from Eq.~\eqref{eq:ch1.maxwell-stress}, this means that scalar summaries of
mechanical loading remain frame-independent even when component entries change.

\begin{figure}[t]
\centering
\fbox{\begin{minipage}{0.92\linewidth}
\centering
\textbf{Observable objectivity test pipeline}\\[2mm]
Input: electromagnetic state \(\mathcal{F}\).\\
Branch 1 (scalar): \(\mathcal{I}[\mathcal{F}] \stackrel{?}{=}
\mathcal{I}[\mathcal{A}_R\mathcal{F}]\).\\
Branch 2 (vector): \(\mathbf{Q}[\mathcal{F}] \stackrel{?}{=} R^{-1}
\mathbf{Q}[\mathcal{A}_R\mathcal{F}]\).\\
Branch 3 (tensor): \(\mathbf{T}\mapsto (I_1,I_2)\) from trace invariants.\\[1mm]
Decision: pass any branch \(\Rightarrow\) frame-objective report;
fail all branches \(\Rightarrow\) coordinate artifact.
\end{minipage}}
\caption{Conceptual screening diagram for frame-independent observables. Each branch corresponds to Eqs.~\eqref{eq:ch1.scalar-observable-invariance}--\eqref{eq:ch1.tensor-principal-invariants}.}
\label{fig:ch1.observable-objectivity-pipeline}
\end{figure}

Figure~\ref{fig:ch1.observable-objectivity-pipeline} converts the algebra into a
usable reporting protocol. The scalar branch tests strict invariance, the vector
branch tests proper rotational transport of direction-valued measurements, and
the tensor branch compresses second-rank content into invariants that are immune
to basis choice. The diagram is not decorative: it encodes the admissibility
filter used before interpreting angular signatures as transport evidence.

For Chapter~\ref{ch:01}, this filter is applied to the directional
angular-momentum map \(\mathbf{\Psi}_J\) from
Eq.~\eqref{eq:ch1.directional-am-flux} and directional active power density
\(\Pi\) from Eq.~\eqref{eq:ch1.directional-power-density}. Define the
dimensionless sector index over \(\Delta\Omega\subseteq 4\pi\):
\begin{equation}
\label{eq:ch1.objective-sector-index}
\Gamma_J(\Delta\Omega;r)=
\frac{\left\|\int_{\Delta\Omega}\mathbf{\Psi}_J(\theta,\phi;r)\,d\Omega\right\|}
{\int_{\Delta\Omega}\Pi(\theta,\phi;r)\,d\Omega+\epsilon_\Pi},
\end{equation}
where \(\|\cdot\|\) is Euclidean norm and \(\epsilon_\Pi>0\) is a small
regularization constant. Under rigid rotation, numerator and denominator are both
reorganized directionally, but \(\Gamma_J\) remains scalar-objective when the
transformed fields satisfy Eqs.~\eqref{eq:ch1.scalar-observable-invariance} and
\eqref{eq:ch1.vector-observable-covariance}. Physically,
\(\Gamma_J\) separates genuine geometry-supported transport strength from
frame-dependent component bookkeeping \cite{harrington2001,collin2001}.

The validity range is restricted to proper rotations in Euclidean space and to
observable constructions built from gauge-robust field quantities. Gauge-dependent
local decompositions and non-Euclidean coordinate manifolds require additional
care that is deferred to Section~\ref{sec:ch01.4}.

Section~\ref{sec:ch01.3} is now complete in substance: covariance machinery,
origin sensitivity, and invariant observables are jointly in place to prevent
coordinate artifacts from entering subsequent gauge discussions.

\paragraph{Section 1.3 closure and forward link.}
After Section~\ref{sec:ch01.3}, rotational covariance and frame-invariant
interpretation are established as mandatory constraints on admissible angular
claims. The chapter next addresses gauge structure in
Section~\ref{sec:ch01.4}, where representational freedom and physical
observables are separated with explicit care.


\section{Gauge Structure and Physical Observables}
\label{sec:ch01.4}

Section~\ref{sec:ch01.4} treats potentials, gauge transformations, and helicity
with emphasis on what survives as an observable quantity. The section is
constructed to prevent local decomposition language from being misread as
gauge-independent physics; every assertion is filtered through an explicit
observability criterion before it is carried forward.


\subsection{Potentials and gauge transformations}
\label{sec:ch01.4.1}
% TOC tag: [HOME]

Potential theory in Maxwell electrodynamics is a representation framework, not an
extra layer of physics: the measurable fields are unchanged under an entire class
of potential transformations. Establishing that class is essential before any
frame-independent observable is interpreted in Section~\ref{sec:ch01.4}.

Let \(\mathbf{A}(\mathbf{r})\) denote vector potential, \(\phi(\mathbf{r})\) scalar
potential, \(\E(\mathbf{r})\) electric phasor field, and \(\mathbf{B}(\mathbf{r})\)
magnetic flux density, with \(j=\sqrt{-1}\), angular frequency \(\omega>0\), and
phasor sign convention fixed for this chapter. Their defining relations are
\begin{equation}
\label{eq:ch1.potential-field-relations}
\mathbf{B}=\nabla\times\mathbf{A},\qquad
\E=-\nabla\phi+j\omega\mathbf{A},
\end{equation}
where \(\nabla\times\) is curl and \(\nabla\) is gradient. The first relation
builds magnetic solenoidality into representation form; the second separates
irrotational electric contribution from inductive contribution. Equation
\eqref{eq:ch1.potential-field-relations} is fully consistent with the previously
established field equations in Eq.~\eqref{eq:ch1.maxwell-local}
\cite{jackson1999,harrington2001}.

Gauge freedom is generated by any sufficiently smooth scalar gauge function
\(\chi(\mathbf{r})\), producing
\begin{equation}
\label{eq:ch1.gauge-transform}
\mathbf{A}'=\mathbf{A}+\nabla\chi,\qquad
\phi'=\phi-j\omega\chi.
\end{equation}
The operand role is explicit: \(\nabla\chi\) alters only longitudinal potential
content, while \(-j\omega\chi\) offsets scalar potential so that reconstructed
\(\E\) and \(\mathbf{B}\) from Eq.~\eqref{eq:ch1.potential-field-relations} remain
unchanged. Physically, Eq.~\eqref{eq:ch1.gauge-transform} identifies an
equivalence orbit of potential descriptions for a single electromagnetic state.

\begin{figure}[t]
\centering
\fbox{\begin{minipage}{0.92\linewidth}
\centering
\textbf{Gauge-orbit and observable pipeline for one field state}\\[2mm]
Layer 1 (representation): nodes \((\mathbf{A},\phi)\), \((\mathbf{A}',\phi')\), \((\mathbf{A}'',\phi'')\).\\
Links on Layer 1: gauge updates by \(\chi_1,\chi_2,\ldots\) from
Eq.~\eqref{eq:ch1.gauge-transform}.\\
Layer 2 (physical): single node \((\E,\mathbf{B})\) obtained from
Eq.~\eqref{eq:ch1.potential-field-relations}.\\
Layer 3 (reporting): scalar observable \(\mathcal{I}\) satisfying
\(\mathcal{I}[\mathbf{A},\phi]=\mathcal{I}[\mathbf{A}',\phi']\).
\end{minipage}}
\caption{Conceptual schematic linking gauge-equivalent potential representations to one physical field state and to a frame/gauge-independent reporting scalar.}
\label{fig:ch1.gauge-orbit}
\end{figure}

Figure~\ref{fig:ch1.gauge-orbit} introduces the observability filter used in this
section: changes on Layer 1 are admissible representation changes; Layer 2 is the
physical content; Layer 3 is an accepted report only if it is invariant across
the gauge orbit. This is the gauge analog of the frame-objectivity discipline
already established in Eq.~\eqref{eq:ch1.scalar-observable-invariance}.

To convert orbit language into a testable criterion, define a scalar reporting
functional \(\mathcal{I}\) on potential pairs. Gauge admissibility requires
\begin{equation}
\label{eq:ch1.gauge-observable-invariance}
\mathcal{I}[\mathbf{A}',\phi']=\mathcal{I}[\mathbf{A},\phi]
\quad\text{for all }\chi\text{ used in Eq.~\eqref{eq:ch1.gauge-transform}}.
\end{equation}
Equation~\eqref{eq:ch1.gauge-observable-invariance} is a chapter-level criterion in
this chapter: any local density that fails it may still be mathematically useful,
but cannot be promoted to gauge-robust observable status.

Representation closure is obtained by imposing Lorenz gauge,
\begin{equation}
\label{eq:ch1.lorenz-gauge-phasor}
\nabla\cdot\mathbf{A}-j\omega\mu\varepsilon\,\phi=0,
\end{equation}
where \(\nabla\cdot\) is divergence, \(\mu\) is permeability, and \(\varepsilon\)
is permittivity. The divergence term suppresses unconstrained longitudinal
freedom, while \(-j\omega\mu\varepsilon\,\phi\) couples scalar potential to
harmonic phase evolution. Mathematically, Eq.~\eqref{eq:ch1.lorenz-gauge-phasor}
regularizes the potential system into a determinate class; physically, it avoids
non-propagative ambiguity in radiation-oriented interpretation
\cite{collin2001,balanis2016}.

The present derivation assumes smooth fields in simply connected regions with
linear constitutive response and single-frequency excitation. Topological gauge
patching, nonlocal constitutive kernels, and strongly nonlinear media require
generalizations outside Chapter~\ref{ch:01}.

With Eqs.~\eqref{eq:ch1.gauge-transform} and
\eqref{eq:ch1.gauge-observable-invariance} fixed, the next subsection can
separate gauge-dependent local densities from quantities admissible as robust
observables.


\subsection{Gauge dependence of local densities}
\label{sec:ch01.4.2}
% TOC tag: [HOME]

Gauge freedom does not threaten measurable fields, yet it can alter local
density formulas that are written directly in terms of potentials. This
subsection establishes a precise diagnostic: which local expressions are genuine
observables and which are representation-dependent bookkeeping terms.

Let \(\mathbf{A}\) and \(\phi\) denote a potential pair and
\((\mathbf{A}',\phi')\) the transformed pair from
Eq.~\eqref{eq:ch1.gauge-transform}. Let \(\mathbf{E}\) and \(\mathbf{B}\) be the
physical fields reconstructed through Eq.~\eqref{eq:ch1.potential-field-relations},
and let \(\mathbf{d}_{\mathrm{loc}}(\mathbf{r})\) be a candidate local density
vector at position \(\mathbf{r}\). The dependence structure can be written as
\begin{equation}
\label{eq:ch1.local-density-functional}
\mathbf{d}_{\mathrm{loc}}(\mathbf{r})
=
\mathcal{D}\!\left[\mathbf{A},\phi,\mathbf{E},\mathbf{B}\right](\mathbf{r}),
\end{equation}
where \(\mathcal{D}\) is a density-construction functional. If
\(\mathcal{D}\) uses potentials explicitly, then the gauge update in
Eq.~\eqref{eq:ch1.gauge-transform} generally induces
\begin{equation}
\label{eq:ch1.local-density-gauge-shift}
\delta_\chi \mathbf{d}_{\mathrm{loc}}
\equiv
\mathbf{d}_{\mathrm{loc}}' - \mathbf{d}_{\mathrm{loc}}
=
\nabla\cdot\mathbf{\Xi}_\chi + \mathbf{\Psi}_\chi,
\end{equation}
where \(\mathbf{\Xi}_\chi\) is a gauge-generated flux tensor and
\(\mathbf{\Psi}_\chi\) is a residual gauge source term. Equation
\eqref{eq:ch1.local-density-gauge-shift} is the core diagnostic of this
subsection: local density objectivity fails whenever either term is nonzero.

The roles of the two terms are distinct. The divergence contribution
\(\nabla\cdot\mathbf{\Xi}_\chi\) can integrate to boundary transport only; the
residual term \(\mathbf{\Psi}_\chi\) changes local value even in closed domains.
Thus, a density may be locally gauge-dependent but still yield gauge-robust
integrated totals under admissible boundary conditions. Let \(V\) be a finite
volume with boundary \(\partial V\) and outward unit normal
\(\hat{\mathbf{n}}\). Then
\begin{equation}
\label{eq:ch1.integrated-density-gauge-test}
\int_V \delta_\chi \mathbf{d}_{\mathrm{loc}}\,dV
=
\int_{\partial V}\mathbf{\Xi}_\chi\cdot\hat{\mathbf{n}}\,dS
\;+\;
\int_V \mathbf{\Psi}_\chi\,dV.
\end{equation}
Equation~\eqref{eq:ch1.integrated-density-gauge-test} follows by divergence
theorem and separates boundary leakage from volumetric gauge contamination.
Mathematically, the left side is gauge variation of integrated density; physically,
it tests whether a reported quantity is robust or representation-contaminated.

\begin{figure}[t]
\centering
\fbox{\begin{minipage}{0.92\linewidth}
\centering
\textbf{Gauge-dependence pathway for local densities}\\[2mm]
Node 1: potential pair \((\mathbf{A},\phi)\).\\
Node 2: transformed pair \((\mathbf{A}',\phi')\) via
Eq.~\eqref{eq:ch1.gauge-transform}.\\
Node 3: local density outputs \(\mathbf{d}_{\mathrm{loc}}\) and
\(\mathbf{d}_{\mathrm{loc}}'\).\\
Difference arrow: \(\delta_\chi \mathbf{d}_{\mathrm{loc}}
=\nabla\cdot\mathbf{\Xi}_\chi+\mathbf{\Psi}_\chi\).\\
Volume box: boundary channel \(\int_{\partial V}\mathbf{\Xi}_\chi\cdot\hat{\mathbf{n}}\,dS\)
plus residual channel \(\int_V\mathbf{\Psi}_\chi\,dV\).
\end{minipage}}
\caption{Conceptual schematic for the gauge variation test in Eqs.~\eqref{eq:ch1.local-density-gauge-shift} and \eqref{eq:ch1.integrated-density-gauge-test}.}
\label{fig:ch1.gauge-local-density-test}
\end{figure}

Figure~\ref{fig:ch1.gauge-local-density-test} visualizes the logical sequence:
gauge-related potential representatives feed a candidate local density functional,
and the resulting difference is decomposed into transport-type and residual
channels. Only when both channels vanish (or cancel under controlled boundary
admissibility) can the density be promoted toward observable status.

The operational criterion that aligns this subsection with
Eq.~\eqref{eq:ch1.gauge-observable-invariance} is
\begin{equation}
\label{eq:ch1.local-observable-admissibility}
\delta_\chi \mathbf{d}_{\mathrm{loc}}=\mathbf{0}
\quad\text{or}\quad
\int_V \delta_\chi \mathbf{d}_{\mathrm{loc}}\,dV=\mathbf{0}
\ \text{under stated boundary conditions}.
\end{equation}
Equation~\eqref{eq:ch1.local-observable-admissibility} does not force every
useful local decomposition to be gauge-invariant pointwise; it states the exact
conditions under which such decompositions may still be interpreted without
violating observable objectivity \cite{jackson1999,collin2001}.

Within Chapter~\ref{ch:01}, this result is used as a filter on local
spin-orbit-style density narratives: forms that fail
Eq.~\eqref{eq:ch1.local-observable-admissibility} remain representation tools,
whereas forms that pass can be linked to measurable channels from
Section~\ref{sec:ch01.5}. This distinction prevents gauge artefacts from being
misreported as physical transport structure \cite{harrington2001,balanis2016}.

The derivation assumes smooth fields, linear media, and boundary regularity
sufficient for divergence theorem use. Singular gauge choices, topological branch
cuts, and nonlocal constitutive operators require generalized distributional
treatment outside the present subsection.

Having separated gauge-dependent local densities from admissible observable
content, the next subsection formalizes gauge-invariant observables alongside
helicity and dual-symmetry structure.


\subsection{Gauge-invariant observables, helicity, and dual symmetry}
\label{sec:ch01.4.3}
% TOC tag: [HOME]

The preceding subsection established that local expressions can be useful yet
still fail observable status when gauge variation survives. The present task is
to identify a class of quantities that remains admissible under both gauge
transformations and rigid-frame changes, and then to connect that class with
helicity transport and electromagnetic dual symmetry.

Consider a linear, isotropic, source-free region with permittivity
\(\varepsilon\), permeability \(\mu\), and intrinsic impedance
\(\eta=\sqrt{\mu/\varepsilon}\). Let \(\E\) and \(\H\) denote phasor electric and
magnetic fields, and let \(\mathcal{F}\) denote the field state. A scalar report
\(\mathcal{O}_s\) and vector report \(\mathbf{O}_v\) are admitted as robust
observables only if they satisfy
\begin{equation}
\label{eq:ch1.observable-admissibility-triad}
\mathcal{O}_s[\mathbf{A}',\phi']=\mathcal{O}_s[\mathbf{A},\phi],\qquad
\mathcal{O}_s[\mathcal{A}_R\mathcal{F}]=\mathcal{O}_s[\mathcal{F}],\qquad
\mathbf{O}_v[\mathcal{A}_R\mathcal{F}]=R\,\mathbf{O}_v[\mathcal{F}],
\end{equation}
where \((\mathbf{A},\phi)\to(\mathbf{A}',\phi')\) follows Eq.~\eqref{eq:ch1.gauge-transform},
\(\mathcal{A}_R\) is the active rotation from
Eq.~\eqref{eq:ch1.active-rotation-action}, and \(R\in\mathrm{SO}(3)\). The first
condition enforces gauge robustness, the second scalar objectivity, and the third
vector covariance; together they specialize the criteria in
Eqs.~\eqref{eq:ch1.gauge-observable-invariance},
\eqref{eq:ch1.scalar-observable-invariance}, and
\eqref{eq:ch1.vector-observable-covariance} into one acceptance test.

Within that admissible class, a gauge-robust local helicity-density functional is
introduced as
\begin{equation}
\label{eq:ch1.helicity-density}
h(\mathbf{r})=
\frac{1}{2\omega}\Im\!\left\{
\varepsilon\,\E^*\cdot(\curl\E)
\;+\;
\mu\,\H^*\cdot(\curl\H)
\right\},
\end{equation}
where \(\Im\{\cdot\}\) extracts imaginary part, \((\cdot)^*\) denotes complex
conjugation, and \(\curl\) is \(\nabla\times\). The first term measures electric
chiral twisting weighted by \(\varepsilon\); the second term measures magnetic
twisting weighted by \(\mu\). Their sum is normalized by \(2\omega\), giving a
frequency-scaled pseudoscalar that changes sign under handedness reversal yet is
unchanged by gauge relabeling of potentials.

Noether structure associated with dual-symmetry rotation yields the local balance
\begin{equation}
\label{eq:ch1.helicity-balance}
\frac{\partial h}{\partial t}
\;+\;
\divg\mathbf{f}_h
\;=\;
-q_h,
\end{equation}
where \(\mathbf{f}_h\) is helicity flux density and \(q_h\) is helicity conversion
rate to matter. In source-free homogeneous domains, \(q_h=0\), so helicity is
transported without local generation or annihilation. This relation plays for
helicity the same structural role that Eq.~\eqref{eq:ch1.am-control-balance}
plays for angular momentum.

Dual symmetry itself is represented by a continuous mixing of \(\E\) and \(\H\):
\begin{equation}
\label{eq:ch1.duality-rotation}
\begin{bmatrix}
\E_\xi\\[1mm]
\eta\H_\xi
\end{bmatrix}
=
\begin{bmatrix}
\cos\xi & \sin\xi\\
-\sin\xi & \cos\xi
\end{bmatrix}
\begin{bmatrix}
\E\\[1mm]
\eta\H
\end{bmatrix},
\end{equation}
with dual angle \(\xi\in\mathbb{R}\). The matrix acts as a rotation in the
\((\E,\eta\H)\) field plane: electric and magnetic components are redistributed
without altering source-free Maxwell structure.

\begin{figure}[t]
\centering
\fbox{\begin{minipage}{0.92\linewidth}
\centering
\textbf{Helicity-duality observable map}\\[2mm]
Left block: admissibility gate from Eq.~\eqref{eq:ch1.observable-admissibility-triad}.\\
Middle block: duality rotation in Eq.~\eqref{eq:ch1.duality-rotation}
acting in the \((\E,\eta\H)\) plane with angle \(\xi\).\\
Right block: helicity channel outputs \(h\), \(\mathbf{f}_h\), and balance
in Eq.~\eqref{eq:ch1.helicity-balance}.\\
Bottom note: only quantities passing the left block may be interpreted as
frame/gauge-robust observables.
\end{minipage}}
\caption{Conceptual schematic linking gauge-frame admissibility, duality rotation, and helicity transport observables.}
\label{fig:ch1.helicity-duality-map}
\end{figure}

Figure~\ref{fig:ch1.helicity-duality-map} supplies the interpretive thread:
duality rotation by itself is a representation-group action, while helicity
diagnostics become physically reportable only after passing the admissibility gate
of Eq.~\eqref{eq:ch1.observable-admissibility-triad}. This prevents a recurring
misstep in advanced discussions, where elegant field decompositions are treated as
observables without gauge/frame screening.

A compact scalar monitor for helicity content is then defined by
\begin{equation}
\label{eq:ch1.normalized-helicity-index}
\mathcal{H}_n(\mathbf{r})=
\frac{h(\mathbf{r})}
{\frac{1}{4}\!\left(\varepsilon\|\E\|^2+\mu\|\H\|^2\right)+\epsilon_h},
\end{equation}
where \(\|\cdot\|\) is Euclidean norm and \(\epsilon_h>0\) is a small
regularization constant. The denominator is local electromagnetic energy density
scale, so \(\mathcal{H}_n\) expresses handedness content relative to local field
strength. As a normalized pseudoscalar built from admissible ingredients, it is
suited for cross-configuration comparison without coordinate-basis bias.

The argument in this subsection assumes linear isotropic constitutive response and source-free
domain conditions required for exact dual symmetry. Material loss, anisotropy,
magnetoelectric coupling, and strong dispersion break or deform the symmetry and
therefore modify Eq.~\eqref{eq:ch1.helicity-balance}.

With gauge-invariant observables and dual-helicity structure now fixed at foundational
level, the chapter can proceed to operational definitions that connect these
quantities to measurable channels.


\subsection{Operational definitions of observables}
\label{sec:ch01.4.4}
% TOC tag: [HOME]

After gauge/frame admissibility is fixed, the next question is operational: by
what measurement map does an admissible field quantity become an observable
number with uncertainty and reproducibility constraints? The present analysis defines
that map in a form that can be used consistently throughout Chapters~1--4.

Let \(\mathcal{F}\) denote an admissible field state and
\(\mathcal{O}[\mathcal{F}]\) a target observable functional that already passes
Eq.~\eqref{eq:ch1.observable-admissibility-triad}. Let
\(\mathcal{M}_\alpha\) denote a measurement channel indexed by apparatus setting
\(\alpha\), and let \(y_\alpha\) be the recorded datum. The operational model is
\begin{equation}
\label{eq:ch1.operational-measurement-model}
y_\alpha=\mathcal{M}_\alpha[\mathcal{F}] + n_\alpha,
\end{equation}
where \(n_\alpha\) is measurement disturbance (noise, drift, calibration residue).
Mathematically, \(\mathcal{M}_\alpha\) maps fields to instrument space; physically,
Eq.~\eqref{eq:ch1.operational-measurement-model} states that observables are never
read directly from symbolic field expressions but through channel-specific
transduction.

An operational definition is accepted only if the inverse map from data to target
observable is stable under admissible perturbations. Introduce estimator
\(\widehat{\mathcal{O}}\) from the data set \(\mathbf{y}\), and define
\begin{equation}
\label{eq:ch1.operational-consistency}
\widehat{\mathcal{O}}=\mathcal{R}[\mathbf{y}],\qquad
\lim_{\|n\|\to 0}\widehat{\mathcal{O}}=\mathcal{O}[\mathcal{F}],
\end{equation}
where \(\mathcal{R}\) is reconstruction rule and \(\|n\|\) is a norm on the noise
vector. The first relation specifies reconstruction architecture; the limit
criterion enforces consistency. Without Eq.~\eqref{eq:ch1.operational-consistency},
an expression may remain mathematically elegant yet experimentally undefined.

Gauge and frame robustness must survive the measurement map itself. If
\(\mathcal{F}'=\mathcal{A}_R\mathcal{F}\) is a rotated state and
\((\mathbf{A}',\phi')\) is a gauge-equivalent potential representative, operational
admissibility requires
\begin{equation}
\label{eq:ch1.channel-covariance}
\mathcal{M}_\alpha[\mathcal{F}']=\mathcal{T}_\alpha(R)\,\mathcal{M}_\alpha[\mathcal{F}],
\qquad
\mathcal{M}_\alpha[\mathbf{A}',\phi']=\mathcal{M}_\alpha[\mathbf{A},\phi],
\end{equation}
where \(\mathcal{T}_\alpha(R)\) is the known channel transformation law under
rotation. The first identity is channel covariance; the second is channel gauge
invariance. These conditions operationalize the abstract criteria established in
Sections~\ref{sec:ch01.3} and \ref{sec:ch01.4}.

\begin{figure}[t]
\centering
\fbox{\begin{minipage}{0.92\linewidth}
\centering
\textbf{Operational observable pipeline}\\[2mm]
Input layer: admissible field state \(\mathcal{F}\).\\
Symmetry gate: pass Eq.~\eqref{eq:ch1.observable-admissibility-triad}.\\
Channel layer: \(\mathcal{F}\xrightarrow{\mathcal{M}_\alpha}y_\alpha\) with
disturbance \(n_\alpha\) from Eq.~\eqref{eq:ch1.operational-measurement-model}.\\
Reconstruction layer: \(\mathbf{y}\xrightarrow{\mathcal{R}}\widehat{\mathcal{O}}\)
subject to Eq.~\eqref{eq:ch1.operational-consistency}.\\
Output layer: accepted observable only if channel covariance/gauge constraints in
Eq.~\eqref{eq:ch1.channel-covariance} hold.
\end{minipage}}
\caption{Conceptual map from admissible field quantity to operationally defined observable.}
\label{fig:ch1.operational-observable-pipeline}
\end{figure}

Figure~\ref{fig:ch1.operational-observable-pipeline} converts the formalism into
an execution sequence for experiments and simulations. The symmetry gate rejects
non-admissible descriptors before instrumentation; the channel layer enforces
physics of sensing; reconstruction produces the reported quantity; the final gate
ensures that reportability is preserved under the same gauge/frame rules used in
theory.

For torque-related channels in Section~\ref{sec:ch01.5}, an immediate operational
example is
\begin{equation}
\label{eq:ch1.operational-torque-channel}
\tau_z^{\mathrm{meas}}
=
\int_{S}\left(\mathbf{r}\times\mathbf{t}\right)\cdot\hat{\mathbf{z}}\,dS
\;+\;\delta\tau_z,
\end{equation}
where \(\mathbf{r}\) is position vector relative to declared origin,
\(\mathbf{t}\) is traction vector from
Eq.~\eqref{eq:ch1.traction-rotation-covariance}, \(\hat{\mathbf{z}}\) is the
measurement axis, \(S\) is sensing surface, and \(\delta\tau_z\) is channel
disturbance. Operand-wise, the surface integral computes deterministic torque
projection; \(\delta\tau_z\) collects non-ideal channel contributions. Physically,
Eq.~\eqref{eq:ch1.operational-torque-channel} makes explicit why origin
declaration and channel calibration are part of observable definition, not
post-processing detail.

This framework assumes calibrated linear channels and weakly nonlinear disturbance
so that \(\mathcal{R}\) is stable. Strongly nonlinear sensor transfer, unknown
cross-coupling, or incomplete channel models require extended inverse methods
beyond Chapter~\ref{ch:01}.

Operational definition is now fixed at foundational level: an observable is a
symmetry-admissible quantity with a stable, channel-explicit reconstruction rule.
The next subsection uses this foundation to delimit where local spin-orbit
decompositions remain interpretive aids rather than direct observables.


\subsection{Limits of local spin-orbit decompositions}
\label{sec:ch01.4.5}
% TOC tag: [HOME]

Local spin-orbit decomposition is mathematically instructive but operationally
fragile: different admissible representation choices can redistribute local
``spin-like'' and ``orbital-like'' densities without changing any measurable total.
The derivation in this subsection establishes the precise boundary between decomposition utility and
observable legitimacy.

Let \(\mathbf{j}_{\mathrm{em}}(\mathbf{r})\) denote electromagnetic
angular-momentum density already fixed in
Eq.~\eqref{eq:ch1.angular-momentum-balance}. A generic local decomposition takes
the form
\begin{equation}
\label{eq:ch1.local-spin-orbit-split}
\mathbf{j}_{\mathrm{em}}
=
\mathbf{l}_{\mathrm{loc}}
\;+\;
\mathbf{s}_{\mathrm{loc}}
\;+\;
\divg\mathbf{C},
\end{equation}
where \(\mathbf{l}_{\mathrm{loc}}\) is orbital-labeled local density,
\(\mathbf{s}_{\mathrm{loc}}\) is spin-labeled local density, \(\mathbf{C}\) is a
second-rank superpotential tensor, and \(\divg\mathbf{C}\) is the divergence
correction generated by decomposition convention. Equation
\eqref{eq:ch1.local-spin-orbit-split} makes clear that label assignment is not
unique at point level: superpotential terms can move content between
\(\mathbf{l}_{\mathrm{loc}}\) and \(\mathbf{s}_{\mathrm{loc}}\) while preserving
\(\mathbf{j}_{\mathrm{em}}\).

Under a gauge update \((\mathbf{A},\phi)\to(\mathbf{A}',\phi')\) from
Eq.~\eqref{eq:ch1.gauge-transform}, the local pair transforms as
\begin{equation}
\label{eq:ch1.spin-orbit-gauge-redistribution}
\mathbf{s}_{\mathrm{loc}}'
=
\mathbf{s}_{\mathrm{loc}}+\divg\mathbf{U}_\chi+\mathbf{v}_\chi,\qquad
\mathbf{l}_{\mathrm{loc}}'
=
\mathbf{l}_{\mathrm{loc}}-\divg\mathbf{U}_\chi-\mathbf{v}_\chi,
\end{equation}
where \(\mathbf{U}_\chi\) is gauge-generated redistribution tensor and
\(\mathbf{v}_\chi\) is residual redistribution vector. The paired signs enforce
that \(\mathbf{s}_{\mathrm{loc}}'+\mathbf{l}_{\mathrm{loc}}'
=\mathbf{s}_{\mathrm{loc}}+\mathbf{l}_{\mathrm{loc}}\), but individual local
components remain gauge-sensitive. In functional terms,
Eq.~\eqref{eq:ch1.spin-orbit-gauge-redistribution} is a redistribution symmetry;
in physical terms, it explains why local spin-orbit maps can disagree across
formulations while total transport predictions remain consistent.

Integrating over a finite volume \(V\) with boundary \(\partial V\) and outward
unit normal \(\hat{\mathbf{n}}\) yields
\begin{equation}
\label{eq:ch1.spin-orbit-volume-test}
\int_V\!\left(\mathbf{s}_{\mathrm{loc}}'+\mathbf{l}_{\mathrm{loc}}'\right)\!dV
=
\int_V\!\left(\mathbf{s}_{\mathrm{loc}}+\mathbf{l}_{\mathrm{loc}}\right)\!dV
=
\int_V\mathbf{j}_{\mathrm{em}}\,dV
\;-\;
\int_{\partial V}\mathbf{C}\cdot\hat{\mathbf{n}}\,dS.
\end{equation}
The first equality shows decomposition invariance of the sum; the second shows
explicit boundary dependence through \(\mathbf{C}\). Therefore local partitions
can be physically compared only after boundary prescription is declared, exactly
as required by the control-balance logic in
Eq.~\eqref{eq:ch1.am-control-balance}.

\begin{figure}[t]
\centering
\fbox{\begin{minipage}{0.92\linewidth}
\centering
\textbf{Spin-orbit local redistribution map}\\[2mm]
Top node: fixed total density \(\mathbf{j}_{\mathrm{em}}\).\\
Middle nodes: candidate local pairs
\((\mathbf{l}_{\mathrm{loc}},\mathbf{s}_{\mathrm{loc}})\) and
\((\mathbf{l}_{\mathrm{loc}}',\mathbf{s}_{\mathrm{loc}}')\).\\
Horizontal arrow: gauge/decomposition redistribution via
Eq.~\eqref{eq:ch1.spin-orbit-gauge-redistribution}.\\
Bottom branch: volume integral in Eq.~\eqref{eq:ch1.spin-orbit-volume-test} with
boundary correction \(\int_{\partial V}\mathbf{C}\cdot\hat{\mathbf{n}}\,dS\).
\end{minipage}}
\caption{Conceptual schematic showing why local spin-orbit partitions are non-unique even when total angular-momentum density is fixed.}
\label{fig:ch1.spin-orbit-limit-map}
\end{figure}

Figure~\ref{fig:ch1.spin-orbit-limit-map} summarizes the limitation pathway:
redistribution acts on local labels, not on the total conserved content. The
operational consequence is immediate: local spin/orbit images are admissible as
diagnostic descriptors, but not as standalone observables unless they pass the
channel and symmetry tests of Eqs.~\eqref{eq:ch1.channel-covariance} and
\eqref{eq:ch1.gauge-observable-invariance}.

To quantify decomposition fragility for a given field realization, define
\begin{equation}
\label{eq:ch1.spin-orbit-nonuniqueness-index}
\kappa_{\mathrm{so}}(V)=
\frac{\int_V\|\mathbf{s}_{\mathrm{loc}}'-\mathbf{s}_{\mathrm{loc}}\|\,dV}
{\int_V\|\mathbf{j}_{\mathrm{em}}\|\,dV+\epsilon_{\mathrm{so}}},
\end{equation}
where \(\|\cdot\|\) is Euclidean norm and \(\epsilon_{\mathrm{so}}>0\) is a small
regularization constant. Small \(\kappa_{\mathrm{so}}\) indicates decomposition
stability for that setup; large \(\kappa_{\mathrm{so}}\) signals strong
representation sensitivity and therefore reduced direct observability of local
spin/orbit attribution \cite{jackson1999,collin2001}.

The derivation assumes smooth fields, linear constitutive response, and boundary
regularity sufficient for divergence manipulations. Singular gauges, strongly
inhomogeneous media, and nonlocal material response can enlarge redistribution
terms and require generalized treatment not developed in Chapter~\ref{ch:01}.

The limit is therefore explicit and constructive: local spin-orbit decomposition
remains a controlled interpretive lens, whereas frame/gauge-admissible totals and
channel-defined measurements remain the primary observables
\cite{harrington2001,balanis2016}.

\paragraph{Section 1.4 summary and continuation.}
Section~\ref{sec:ch01.4} delineates the boundary between gauge-dependent
descriptors and gauge-robust observables. On that basis, the chapter proceeds to
Section~\ref{sec:ch01.5}, where measurement channels and interpretation limits
are formalized without reintroducing gauge ambiguity.


\section{Measurement and Interpretation}
\label{sec:ch01.5}

Section~\ref{sec:ch01.5} translates formal field statements into experimentally
resolvable quantities. Its purpose is to define how torque, flux, and modal
projections are interpreted under finite apparatus and finite information, while
maintaining strict separation between direct observables and inferred
reconstructions.


\subsection{Field-matter interaction}
\label{sec:ch01.5.1}
% TOC tag: [HOME]

Measurement enters electromagnetics through interaction channels in which field
variables exchange power, momentum, or angular momentum with material degrees of
freedom. A precise interaction map is established here at foundational level so that
later inference statements in Chapter~\ref{ch:01} are judged against explicit
transfer equations rather than against representation-level descriptors.

Let \(V_m\subset\Omega\) denote the material interaction region, with boundary
\(\partial V_m\), outward unit normal \(\hat{\mathbf{n}}\), and volume element
\(dV\). Let \(\E(\mathbf{r})\) and \(\H(\mathbf{r})\) be phasor fields under the
\(\exp(-j\omega t)\) convention, let \(\mathbf{J}_c(\mathbf{r})\) be conduction-current
density inside matter, and let \(\rho(\mathbf{r})\) be free-charge density. The
time-averaged electromagnetic power-transfer density into matter is
\begin{equation}
\label{eq:ch1.material-power-density}
p_m(\mathbf{r})=\frac{1}{2}\Re\!\left\{\E\cdot\mathbf{J}_c^*\right\},
\end{equation}
where \((\cdot)^*\) denotes complex conjugation and \(\Re\{\cdot\}\) extracts the
real part. Operand by operand, \(\E\cdot\mathbf{J}_c^*\) is local field-current work-rate
density in complex form, and the factor \(1/2\) converts phasor bilinear form to
time average. As a complete relation,
Eq.~\eqref{eq:ch1.material-power-density} identifies the local channel by which
electromagnetic energy is converted into measurable material response such as
heating or mechanical actuation \cite{jackson1999,collin2001}.

Mechanical transfer is described by the cycle-averaged Lorentz-force density
\begin{equation}
\label{eq:ch1.material-force-density}
\mathbf{f}_m(\mathbf{r})
=\frac{1}{2}\Re\!\left\{
\rho\,\E+\mathbf{J}_c\times\mu\,\H^*
\right\},
\end{equation}
where \(\mu\) is permeability and \(\times\) is vector cross product. The first
term \(\rho\,\E\) is electric-force density on free charge, while
\(\mathbf{J}_c\times\mu\H^*\) is magnetic-force density on moving charge carriers. Taken
together, Eq.~\eqref{eq:ch1.material-force-density} is the local force-transfer
law linking measurable mechanical loading to field state and material current.

Integrating over \(V_m\) gives net power and net force channels:
\[
P_m=\int_{V_m}p_m\,dV,\qquad
\mathbf{F}_m=\int_{V_m}\mathbf{f}_m\,dV.
\]
These unnumbered expressions are direct integral uses of
Eqs.~\eqref{eq:ch1.material-power-density}--\eqref{eq:ch1.material-force-density}
and therefore do not introduce new theorem-level content. Their physical role is
operational: they define what a calorimetric or force-sensing apparatus should
recover in the absence of channel distortion.

\begin{figure}[t]
\centering
\fbox{\begin{minipage}{0.92\linewidth}
\centering
\textbf{Field-to-matter transfer schematic}\\[2mm]
Input block: field state \((\E,\H)\) over interaction region \(V_m\).\\
Channel 1: \(p_m=\frac{1}{2}\Re\{\E\cdot\mathbf{J}_c^*\}\) integrated to \(P_m\).\\
Channel 2: \(\mathbf{f}_m=\frac{1}{2}\Re\{\rho\E+\mathbf{J}_c\times\mu\H^*\}\)
integrated to \(\mathbf{F}_m\).\\
Output block: measurable thermal/mechanical response under stated boundary and
material assumptions.
\end{minipage}}
\caption{Conceptual map of local-to-global field-matter transfer channels in Subsection~\ref{sec:ch01.5.1}.}
\label{fig:ch1.field-matter-transfer}
\end{figure}

Figure~\ref{fig:ch1.field-matter-transfer} clarifies the operational structure:
measurement does not access abstract field amplitudes directly, but accesses
material response after transfer through power and force channels. This is why a
claim about ``observed angular structure'' remains incomplete unless the channel
from fields to matter is specified.

The validity envelope of this subsection is limited to linear local constitutive
response, harmonic steady state, and continuum averaging scales for charge and
current densities. Strongly nonlinear transport, discrete-particle regimes, or
ultrafast transients require generalized models and are outside the present foundational
establishment \cite{harrington2001,balanis2016}.

The next subsection converts the stress-tensor framework from
Subsection~\ref{sec:ch01.1.4} into explicit torque and momentum measurement
relations.


\subsection{Torque, momentum flux, and stress measurements}
\label{sec:ch01.5.2}
% TOC tag: [USE SS1.1.4]

Stress-based measurement in this chapter is inherited directly from the owner
framework in Subsection~\ref{sec:ch01.1.4}; no re-derivation of field momentum
balance is introduced here. The objective is strictly operational: map the
already-established stress tensor to instrument outputs for force and torque.

Let \(\mathbf{T}\) be Maxwell stress tensor from
Eq.~\eqref{eq:ch1.maxwell-stress}, let \(S_m\) be the sensing surface around a
test body, and let \(\hat{\mathbf{n}}\) denote outward surface normal on \(S_m\).
The traction density is \(\mathbf{t}=\mathbf{T}\cdot\hat{\mathbf{n}}\). For a
declared reference origin \(\mathbf{r}_0\), measured force vector
\(\mathbf{y}_F\), and measured torque component \(y_\tau\) about axis unit vector
\(\hat{\mathbf{u}}_\tau\), a channel-calibrated measurement model is
\begin{equation}
\label{eq:ch1.stress-measurement-model}
\mathbf{y}_F
=\mathbf{K}_F\!\left(\int_{S_m}\mathbf{T}\cdot\hat{\mathbf{n}}\,dS\right)
+\mathbf{n}_F,\qquad
y_\tau
=K_\tau\!\left(
\int_{S_m}\!\big[(\mathbf{r}-\mathbf{r}_0)\times(\mathbf{T}\cdot\hat{\mathbf{n}})\big]
\cdot\hat{\mathbf{u}}_\tau\,dS
\right)+n_\tau,
\end{equation}
where \(\mathbf{K}_F\) and \(K_\tau\) are calibration operators/scalars, and
\(\mathbf{n}_F,n_\tau\) collect channel disturbances. The two surface integrals are
deterministic electromagnetic load predictions inherited from
Eq.~\eqref{eq:ch1.maxwell-stress}; calibration and disturbance terms convert that
prediction into practical measurement channels.

At mathematical level, Eq.~\eqref{eq:ch1.stress-measurement-model} composes
field-to-traction mapping, surface integration, and instrument transfer in one
statement. At physical level, it explains why two laboratories can report
different torque values from the same field realization unless origin declaration,
axis definition, and calibration normalization are matched.

\begin{figure}[t]
\centering
\fbox{\begin{minipage}{0.92\linewidth}
\centering
\textbf{Stress-to-instrument mapping for torque channels}\\[2mm]
Surface \(S_m\) encloses the test body; traction
\(\mathbf{t}=\mathbf{T}\cdot\hat{\mathbf{n}}\) is integrated for force and moment.\\
Reference origin \(\mathbf{r}_0\) and axis \(\hat{\mathbf{u}}_\tau\) define the
reported torque component.\\
Calibration blocks \(\mathbf{K}_F,K_\tau\) and disturbance inputs
\(\mathbf{n}_F,n_\tau\) complete the measured outputs
\((\mathbf{y}_F,y_\tau)\).
\end{minipage}}
\caption{Schematic of stress-flux integration and calibration stages used in Subsection~\ref{sec:ch01.5.2}.}
\label{fig:ch1.stress-measurement-map}
\end{figure}

Figure~\ref{fig:ch1.stress-measurement-map} is the practical counterpart of
Subsection~\ref{sec:ch01.1.4}: the stress tensor remains the field-theoretic
object, while this subsection fixes only how that object is sampled and reported.
This reference-only inheritance discipline prevents recursion while preserving experimental
readability.

The model assumes quasi-static sensor response over one cycle-average window and
surface geometry known to sufficient precision for traction integration. Large
vibration bandwidth, unknown moving boundaries, or unresolved near-contact fields
can invalidate the simple calibration form in
Eq.~\eqref{eq:ch1.stress-measurement-model} and require dynamic inverse
identification \cite{jackson1999,harrington2001}.

The section next states how modal coefficients are estimated from finite
measurements and where those estimates become ill-conditioned.


\subsection{Modal projections and measurement limits}
\label{sec:ch01.5.3}
% TOC tag: [HOME]

Angular interpretation often proceeds through modal coefficients, but those
coefficients are not primitive observables; they are reconstructions from finite
channel data. The projection model introduced below exposes the resolution
limits imposed by finite aperture, finite bandwidth, and noise.

Let \(\Gamma_m\) denote a measurement surface, let
\(\{\mathbf{w}_q(\mathbf{r})\}_{q=1}^{Q}\) be known channel weighting functions on
\(\Gamma_m\), and let \(y_q\) be measured channel data. For tangential electric
field \(\E_t\), channel outputs satisfy
\begin{equation}
\label{eq:ch1.modal-projection-data}
y_q=\int_{\Gamma_m}\mathbf{w}_q^*(\mathbf{r})\cdot\E_t(\mathbf{r})\,dS+n_q,
\qquad q=1,\dots,Q,
\end{equation}
where \(n_q\) is measurement noise in channel \(q\). Each integral is a weighted
projection of \(\E_t\) onto channel kernel \(\mathbf{w}_q\), so the data vector
\(\mathbf{y}\) captures only the subspace interrogated by the chosen sensing set.

Represent \(\E_t\) in a model basis \(\{\boldsymbol{\psi}_m\}_{m=1}^{M}\) with
coefficients \(a_m\): \(\E_t\approx\sum_{m=1}^{M}a_m\boldsymbol{\psi}_m\). Then
Eq.~\eqref{eq:ch1.modal-projection-data} becomes the finite-dimensional inverse
system \(\mathbf{y}=G\mathbf{a}+\mathbf{n}\), where
\(G_{qm}=\int_{\Gamma_m}\mathbf{w}_q^*\cdot\boldsymbol{\psi}_m\,dS\). A stable
regularized estimator is
\begin{equation}
\label{eq:ch1.modal-regularized-estimator}
\widehat{\mathbf{a}}_\lambda
=\left(G^*G+\lambda I\right)^{-1}G^*\mathbf{y},
\end{equation}
where \(G^*\) is conjugate transpose, \(I\) is identity matrix, and \(\lambda>0\)
is regularization parameter. Operator-wise, \(G^*G\) encodes channel sensitivity
to basis directions, \(\lambda I\) suppresses small-singular-value amplification,
and left multiplication by \((G^*G+\lambda I)^{-1}G^*\) yields the reconstructed
coefficient vector.

As a complete inverse relation,
Eq.~\eqref{eq:ch1.modal-regularized-estimator} formalizes a central measurement
fact: modal descriptions are model-dependent estimates whose reliability depends
on channel geometry and conditioning, not on phase pattern alone. This condition
is the section-level safeguard against over-interpreting visually appealing but
weakly observable modal decompositions \cite{collin2001,balanis2016}.

Quantitative resolution is bounded by singular-spectrum decay of \(G\). Let
\(\sigma_{\min}^{(\mathrm{eff})}\) be the smallest retained singular value after
truncation. Then, for perturbation \(\delta\mathbf{y}\), one obtains the
conditioning bound
\[
\|\delta\widehat{\mathbf{a}}_\lambda\|_2
\lesssim
\frac{1}{\sigma_{\min}^{(\mathrm{eff})}+\lambda}\,\|\delta\mathbf{y}\|_2.
\]
This unnumbered inequality is a direct stability reading of
Eq.~\eqref{eq:ch1.modal-regularized-estimator}: weakly observed modal directions
amplify uncertainty unless regularization or acquisition geometry is revised.

\begin{figure}[t]
\centering
\fbox{\begin{minipage}{0.92\linewidth}
\centering
\textbf{Projection-to-estimation pipeline for modal reports}\\[2mm]
Field on \(\Gamma_m\) \(\rightarrow\) weighted projections
\(\{y_q\}\) via Eq.~\eqref{eq:ch1.modal-projection-data}.\\
Data vector \(\mathbf{y}\) \(\rightarrow\) inverse block with matrix \(G\),
regularization \(\lambda\), estimator
\(\widehat{\mathbf{a}}_\lambda\) in Eq.~\eqref{eq:ch1.modal-regularized-estimator}.\\
Conditioning branch: singular-value decay controls reliable modal dimension.
\end{minipage}}
\caption{Conceptual schematic separating measured channel space from reconstructed modal-coefficient space in Subsection~\ref{sec:ch01.5.3}.}
\label{fig:ch1.modal-projection-pipeline}
\end{figure}

Figure~\ref{fig:ch1.modal-projection-pipeline} emphasizes the interpretive
boundary required in this chapter: channels are measured, coefficients are
inferred. The distinction is indispensable when angular labels are compared
across different apertures or sensing protocols.

The present derivation assumes linear measurement kernels, known basis functions,
and perturbations small enough for first-order conditioning analysis. Model
mismatch, nonlinear detectors, or incomplete channel characterization require
broader uncertainty quantification than addressed here.

The final subsection formulates a decision criterion that separates direct
observables from quantities that remain inference-dependent descriptors.


\subsection{Observables versus inferred quantities}
\label{sec:ch01.5.4}
% TOC tag: [HOME]

The section closes by drawing a strict boundary between directly measured
quantities and reconstructed descriptors. This boundary is indispensable for
Chapter~\ref{ch:01}: without it, representation-level elegance can be mistaken
for measurable electromagnetic content.

Let \(\mathcal{O}_d\) denote a direct observable generated by calibrated channel
map \(\mathcal{M}\), and let \(\mathcal{Q}_i\) denote an inferred quantity
computed by estimator \(\mathcal{R}\) from data \(\mathbf{y}\). Using
Eq.~\eqref{eq:ch1.operational-measurement-model}, write
\[
\mathcal{O}_d=\mathcal{M}[\mathcal{F}],\qquad
\mathcal{Q}_i=\mathcal{R}[\mathbf{y}],\qquad
\mathbf{y}=\mathcal{M}[\mathcal{F}]+\mathbf{n}.
\]
These unnumbered relations define roles only: \(\mathcal{O}_d\) belongs to the
channel output space, while \(\mathcal{Q}_i\) belongs to reconstruction space.

A quantity is admitted as operationally robust only if it satisfies symmetry
admissibility from Eq.~\eqref{eq:ch1.channel-covariance} together with bounded
uncertainty amplification. Introduce effective uncertainty gain
\(\kappa_{\mathrm{eff}}\) and direct-channel uncertainty floor
\(\sigma_d\), and define
\begin{equation}
\label{eq:ch1.observable-inference-criterion}
\mathfrak{A}(\mathcal{Q}_i)
=
\frac{\mathrm{Std}[\mathcal{Q}_i]}
{\sigma_d}
\cdot
\kappa_{\mathrm{eff}}.
\end{equation}
Operand-level meaning is explicit: \(\mathrm{Std}[\mathcal{Q}_i]\) is estimator
uncertainty, \(\sigma_d\) is measurement-scale uncertainty of direct channels,
and \(\kappa_{\mathrm{eff}}\) captures inversion sensitivity. As a complete
criterion, Eq.~\eqref{eq:ch1.observable-inference-criterion} quantifies how far
an inferred descriptor departs from direct-observable reliability.

An interpretation decision is then stated by
\begin{equation}
\label{eq:ch1.observable-decision-rule}
\begin{aligned}
\mathfrak{A}(\mathcal{Q}_i)\le \tau_A
&\Longrightarrow
\text{report as constrained observable surrogate},\\
\mathfrak{A}(\mathcal{Q}_i)>\tau_A
&\Longrightarrow
\text{report as inference-only descriptor}.
\end{aligned}
\end{equation}
where \(\tau_A>0\) is an application-specific admissibility threshold fixed by
protocol before analysis. The mathematical function of
Eq.~\eqref{eq:ch1.observable-decision-rule} is classification under explicit
uncertainty control; the physical function is to prevent modal or phase labels
from being over-claimed as directly measured transport quantities.

\begin{figure}[t]
\centering
\fbox{\begin{minipage}{0.92\linewidth}
\centering
\textbf{Observable vs inference decision architecture}\\[2mm]
Measured data path: \(\mathcal{F}\xrightarrow{\mathcal{M}}\mathbf{y}\).\\
Direct report branch: \(\mathbf{y}\rightarrow \mathcal{O}_d\).\\
Inference branch: \(\mathbf{y}\xrightarrow{\mathcal{R}}\mathcal{Q}_i\rightarrow
\mathfrak{A}(\mathcal{Q}_i)\) via Eq.~\eqref{eq:ch1.observable-inference-criterion}.\\
Decision gate: Eq.~\eqref{eq:ch1.observable-decision-rule} assigns
observable-surrogate or inference-only status.
\end{minipage}}
\caption{Conceptual classifier separating direct observables from inference-dependent quantities in Subsection~\ref{sec:ch01.5.4}.}
\label{fig:ch1.observable-inference-classifier}
\end{figure}

Figure~\ref{fig:ch1.observable-inference-classifier} provides the reporting logic
required for the rest of Volume~I: every quoted angular descriptor must carry its
classification tag, so that experimentally anchored quantities are not mixed with
model-contingent reconstructions.

The criterion assumes calibrated channels, declared thresholds, and uncertainty
estimates consistent with the noise model of
Eq.~\eqref{eq:ch1.operational-measurement-model}. In adversarial noise, strong
model mismatch, or nonstationary acquisition, thresholding alone is insufficient
and should be replaced by full Bayesian or worst-case robust inference protocols
\cite{harrington2001,collin2001}.

Section~\ref{sec:ch01.5} now closes with a complete measurement interpretation
stack: field-matter transfer channels, stress-based torque measurement, modal
projection limits, and an explicit observable-versus-inference decision rule.
This stack is the acceptance filter applied to all worked cases in
Section~\ref{sec:ch01.6}.

\paragraph{Section 1.5 consolidation and handoff.}
The measurement architecture established in Section~\ref{sec:ch01.5} converts
field-theoretic statements into channel-explicit reporting rules with declared
uncertainty boundaries. Section~\ref{sec:ch01.6} now stress-tests those rules on
worked examples and counterexamples, so that conservation-consistent angular
claims can be separated from phase-only or inversion-fragile interpretations.


\section{Examples and Counterexamples}
\label{sec:ch01.6}
% TOC tag: [FIG+PROB]

Section~\ref{sec:ch01.6} is the verification layer of Chapter~\ref{ch:01}. Each
subsection is drafted in full problem form, with schematic support and explicit
mathematical steps, to test the theoretical criteria established earlier against
canonical and limiting cases.


\subsection{Angular momentum conservation in radiation}
\label{sec:ch01.6.1}
% TOC tag: [FIG+PROB][USE SS1.1.5]

Consider a finite radiator enclosed by a closed control surface \(S_R\) in the
radiative zone, with outward unit normal \(\hat{\mathbf{n}}\), position vector
\(\mathbf{r}\) measured from declared origin \(\mathbf{r}_0\), and time-averaged
Maxwell stress tensor \(\mathbf{T}\) inherited from
Subsection~\ref{sec:ch01.1.4}. The problem is to predict measurable mechanical
torque about axis \(\hat{\mathbf{z}}\) from radiated power and angular mode
content, without re-deriving the conservation law already fixed in
Subsection~\ref{sec:ch01.1.5}.

For steady harmonic operation, storage variation over one cycle-average window is
zero, so the angular-momentum balance from
Eq.~\eqref{eq:ch1.angular-momentum-balance} reduces to a surface-flux/torque
equation:
\begin{equation}
\label{eq:ch1.ex61-torque-flux}
\tau_z^{\mathrm{mech}}
=
-\int_{S_R}
\big[\mathbf{r}\times(\mathbf{T}\cdot\hat{\mathbf{n}})\big]
\cdot\hat{\mathbf{z}}\,dS,
\end{equation}
where \(\tau_z^{\mathrm{mech}}\) is measured mechanical torque component on the
load. The integrand is local angular-momentum flux density crossing \(S_R\),
projected on \(\hat{\mathbf{z}}\); the surface integral is total exported
electromagnetic angular momentum per unit time.

Let \(P_{\mathrm{rad}}\) denote radiated power through \(S_R\), and define
dimensionless angular transport factor \(\chi_J\) by
\[
\chi_J=\frac{\dot J_z}{P_{\mathrm{rad}}},\qquad
\dot J_z=
\int_{S_R}
\big[\mathbf{r}\times(\mathbf{T}\cdot\hat{\mathbf{n}})\big]
\cdot\hat{\mathbf{z}}\,dS.
\]
In the common single-helical-channel limit with azimuthal index \(m\in\mathbb{Z}\)
and circular-polarization helicity \(\sigma\in\{-1,+1\}\), the channel factor is
\(\chi_J\approx (m+\sigma)/\omega\), where \(\omega\) is angular frequency.
Substituting into Eq.~\eqref{eq:ch1.ex61-torque-flux} yields the application law
\begin{equation}
\label{eq:ch1.ex61-torque-power}
\tau_z^{\mathrm{mech}}
\approx
-\frac{m+\sigma}{\omega}\,P_{\mathrm{rad}}.
\end{equation}
Equation~\eqref{eq:ch1.ex61-torque-power} is not a new conservation theorem; it
is a regime-specific evaluator obtained by applying the conservation
statement of Subsection~\ref{sec:ch01.1.5} to a radiative control surface.

\begin{figure}[t]
\centering
\fbox{\begin{minipage}{0.92\linewidth}
\centering
\textbf{Problem geometry: radiative angular-momentum flux to torque}\\[2mm]
Central radiator at origin \(\mathbf{r}_0\), enclosed by far-zone control surface
\(S_R\).\\
Outward flux channel:
\(\big[\mathbf{r}\times(\mathbf{T}\cdot\hat{\mathbf{n}})\big]\cdot\hat{\mathbf{z}}\).\\
Mechanical load on axis \(\hat{\mathbf{z}}\) reports
\(\tau_z^{\mathrm{mech}}\).\\
Power monitor on \(S_R\) reports \(P_{\mathrm{rad}}\).
\end{minipage}}
\caption{Worked-problem schematic for Subsection~\ref{sec:ch01.6.1}: mapping radiative angular-momentum flux to measurable torque.}
\label{fig:ch1.ex61-radiation-torque}
\end{figure}

Figure~\ref{fig:ch1.ex61-radiation-torque} expresses the full measurement chain:
field solution \(\rightarrow\) stress-flux evaluation \(\rightarrow\) torque
readout. The physical interpretation is immediate: once conservation holds, a
nonzero exported angular-momentum flux must appear as opposite-sign torque on the
radiator-load system. The measured torque therefore tests transport reality, not
phase aesthetics \cite{jackson1999,harrington2001}.

The result is valid when \(S_R\) lies in a region where radiative flux dominates
reactive recirculation, material back-action on the source is weak over one cycle
average, and modal leakage to uncontrolled channels is negligible. Strong
near-field storage, multimode cross-coupling, or moving boundaries require a
broader control model \cite{collin2001,jackson1999}.

The next problem deliberately constructs a field with conspicuous phase structure
but vanishing net angular-momentum transport, thereby testing necessity versus
sufficiency of phase winding.

\subsection{Phase structure without angular-momentum transport}
\label{sec:ch01.6.2}
% TOC tag: [FIG+PROB][HOME]

Let a monochromatic transverse field on observation surface \(S_o\) be formed by
superposing two equal-amplitude counter-rotating helical channels of the same
radial envelope \(A(\rho)\), cylindrical radius \(\rho\), azimuth \(\phi\), and
integer order \(m>0\):
\begin{equation}
\label{eq:ch1.ex62-counter-rotating-superposition}
E_t(\rho,\phi)
=
A(\rho)e^{j m\phi}
+
A(\rho)e^{-j m\phi}
=
2A(\rho)\cos(m\phi).
\end{equation}
Equation~\eqref{eq:ch1.ex62-counter-rotating-superposition} exhibits pronounced
angular patterning through \(\cos(m\phi)\), including alternating lobes and nodal
lines, so the field is phase-structured by any geometric inspection criterion.

Define modal intensities \(I_+(\rho)=|A_+(\rho)|^2\) and
\(I_-(\rho)=|A_-(\rho)|^2\) for \(+m\) and \(-m\) channels, and let
\(\dot J_z\) denote net angular-momentum transport rate through \(S_o\). For this
paired channel family, the transport balance is
\begin{equation}
\label{eq:ch1.ex62-net-am-transport}
\dot J_z
=
\frac{m}{\omega}
\int_{S_o}\!\big(I_+-I_-\big)\,dS.
\end{equation}
Operand interpretation is explicit: \(m/\omega\) is angular-momentum-per-energy
scale for each channel, while \(I_+-I_-\) is signed channel imbalance. The full
equation states that net transport depends on imbalance, not on phase pattern
visibility.

Under equal excitation \(I_+=I_-\), Eq.~\eqref{eq:ch1.ex62-net-am-transport}
gives
\[
\dot J_z=0
\quad\text{while}\quad
E_t(\rho,\phi)=2A(\rho)\cos(m\phi)\neq 0.
\]
Hence phase structure is not sufficient for angular-momentum transport. This
subsection establishes that insufficiency at foundational level and fixes a critical
interpretation rule used throughout later chapters.

\begin{figure}[t]
\centering
\fbox{\begin{minipage}{0.92\linewidth}
\centering
\textbf{Counterexample geometry: visible phase structure, zero net transport}\\[2mm]
Channel \(+m\): azimuthal circulation in positive sense.\\
Channel \(-m\): azimuthal circulation in negative sense with equal intensity.\\
Superposed field: lobe/nodal pattern in \(\phi\), but channel imbalance is zero,
so net \(\dot J_z\) through \(S_o\) is zero.
\end{minipage}}
\caption{Schematic for Subsection~\ref{sec:ch01.6.2}: phase-pattern richness without net angular-momentum carriage.}
\label{fig:ch1.ex62-phase-no-transport}
\end{figure}

Figure~\ref{fig:ch1.ex62-phase-no-transport} converts the algebra into a
diagnostic statement: visual twisting or lobe count cannot by itself certify
mechanical torque capability. A transport claim requires signed channel imbalance
or equivalent flux criterion, exactly as enforced by
Eq.~\eqref{eq:ch1.ex62-net-am-transport} \cite{jackson1999,collin2001}.

The derivation presumes linear superposition, coherent channel control, and
observation geometry that resolves channel intensities without bias. Incoherent
mixing, polarization-selective loss, or aperture truncation can break the exact
cancelation and produce residual transport.

With this counterexample now fixed, the next subsection tests an additional
failure mode: gauge-sensitive local density narratives that survive algebraic
manipulation but fail observability discipline.

\subsection{Gauge dependence of local angular-momentum densities}
\label{sec:ch01.6.3}
% TOC tag: [FIG+PROB][USE SS1.4.2]

Take a finite volume \(V_g\) with boundary \(\partial V_g\), outward unit normal
\(\hat{\mathbf{n}}\), and two gauge-related potential representatives
\((\mathbf{A},\phi)\) and \((\mathbf{A}',\phi')\). The gauge-variation structure already
established in Subsection~\ref{sec:ch01.4.2} is used here directly; no new gauge
theorem is derived. The problem is to quantify when local angular-momentum
density maps disagree while integrated transport observables remain unchanged.

Let \(\mathbf{j}_{\mathrm{loc}}\) be a candidate local angular-momentum density
constructed from potentials and fields, and define gauge difference
\(\delta_\chi\mathbf{j}_{\mathrm{loc}}
=\mathbf{j}_{\mathrm{loc}}'-\mathbf{j}_{\mathrm{loc}}\).
Applying the inherited structure from Subsection~\ref{sec:ch01.4.2} gives
\begin{equation}
\label{eq:ch1.ex63-gauge-density-balance}
\int_{V_g}\delta_\chi\mathbf{j}_{\mathrm{loc}}\,dV
=
\int_{\partial V_g}\mathbf{\Xi}_\chi\cdot\hat{\mathbf{n}}\,dS
+
\int_{V_g}\mathbf{\Psi}_\chi\,dV,
\end{equation}
where \(\mathbf{\Xi}_\chi\) is gauge-induced flux tensor and \(\mathbf{\Psi}_\chi\)
is residual gauge source term (notations inherited from
Subsection~\ref{sec:ch01.4.2}). The first integral on the right is boundary
redistribution, while the second is volumetric contamination.

Under boundary/material conditions for which
\(\int_{\partial V_g}\mathbf{\Xi}_\chi\cdot\hat{\mathbf{n}}\,dS=0\) and
\(\int_{V_g}\mathbf{\Psi}_\chi\,dV=0\), Eq.~\eqref{eq:ch1.ex63-gauge-density-balance}
reduces to
\[
\int_{V_g}\mathbf{j}_{\mathrm{loc}}'\,dV
=
\int_{V_g}\mathbf{j}_{\mathrm{loc}}\,dV,
\]
even though pointwise maps \(\mathbf{j}_{\mathrm{loc}}'(\mathbf{r})\) and
\(\mathbf{j}_{\mathrm{loc}}(\mathbf{r})\) may differ strongly. Therefore local
gauge-sensitive pictures are admissible only as intermediate descriptors, not as
standalone observables.

\begin{figure}[t]
\centering
\fbox{\begin{minipage}{0.92\linewidth}
\centering
\textbf{Gauge-dependent local map versus gauge-robust integral report}\\[2mm]
Branch 1: \((\mathbf{A},\phi)\rightarrow \mathbf{j}_{\mathrm{loc}}\).\\
Branch 2: \((\mathbf{A}',\phi')\rightarrow \mathbf{j}_{\mathrm{loc}}'\).\\
Difference channel:
\(\delta_\chi\mathbf{j}_{\mathrm{loc}}\rightarrow
\int_{\partial V_g}\mathbf{\Xi}_\chi\cdot\hat{\mathbf{n}}\,dS
+\int_{V_g}\mathbf{\Psi}_\chi\,dV\).\\
Integral report node: equality/non-equality of
\(\int_{V_g}\mathbf{j}_{\mathrm{loc}}\,dV\) across gauges.
\end{minipage}}
\caption{Worked-problem schematic for Subsection~\ref{sec:ch01.6.3}: local gauge sensitivity versus integral observability.}
\label{fig:ch1.ex63-gauge-density}
\end{figure}

Figure~\ref{fig:ch1.ex63-gauge-density} makes the interpretation criterion
operational: a local map may change under gauge update without violating physics,
provided integral observables survive the gauge test. The subsection thus
reinforces non-recursive inheritance discipline by applying, rather than restating, the gauge
admissibility framework of Subsection~\ref{sec:ch01.4.2}
\cite{jackson1999,harrington2001}.

The argument holds when potentials are sufficiently regular for volume/surface
operations and when boundary conditions are declared explicitly. Topological gauge
patching, singular gauges, or unresolved boundary currents can invalidate the
simplified cancellation path.

The final worked case compares reactive and radiative regimes directly, using the
regime-separation principles of Subsection~\ref{sec:ch01.2.4} in measurement form.

\subsection{Reactive versus radiative comparison}
\label{sec:ch01.6.4}
% TOC tag: [FIG+PROB][USE SS1.2.4]

Let \(S_{r_1}\) and \(S_{r_2}\) be concentric spherical observation surfaces,
with radii \(r_1<r_2\), centered on the same source. The objective is to compare
angular-structure interpretation in a reactive-dominant surface \(S_{r_1}\) and a
radiative-dominant surface \(S_{r_2}\), using the regime logic inherited from
Subsection~\ref{sec:ch01.2.4}.

Define active-power and reactive-power surface functionals
\begin{equation}
\label{eq:ch1.ex64-power-functionals}
P(r)=\int_{S_r}\Re\!\left\{\mathbf{S}\cdot\hat{\mathbf{n}}\right\}dS,\qquad
Q(r)=\int_{S_r}\Im\!\left\{\mathbf{S}\cdot\hat{\mathbf{n}}\right\}dS,
\end{equation}
where \(\mathbf{S}\) is complex Poynting vector and \(\Im\{\cdot\}\) extracts
imaginary part. The first integral measures net exported power; the second
measures reactive circulation associated with local storage exchange.

For a torque channel about \(\hat{\mathbf{z}}\), define
\begin{equation}
\label{eq:ch1.ex64-torque-contrast}
\mathcal{C}_\tau
=
\frac{|\tau_z(r_2)|}{|\tau_z(r_1)|+\epsilon_\tau},
\qquad
\epsilon_\tau>0.
\end{equation}
The numerator and denominator are measured torque magnitudes on outer and inner
surfaces; \(\epsilon_\tau\) is a small regularization constant preventing division
singularity. When \(Q(r_1)\gg P(r_1)\) and \(P(r_2)\gg |Q(r_2)|\), one typically
obtains \(\mathcal{C}_\tau\gg 1\): near-zone angular structure is visually rich
yet weakly transport-capable, while far-zone structure carries measurable torque.

\begin{figure}[t]
\centering
\fbox{\begin{minipage}{0.92\linewidth}
\centering
\textbf{Two-surface regime comparison problem}\\[2mm]
Inner surface \(S_{r_1}\): reactive-dominant (\(|Q| \gg P\)).\\
Outer surface \(S_{r_2}\): radiative-dominant (\(P \gg |Q|\)).\\
Both surfaces evaluate \(P(r)\), \(Q(r)\), and torque channel \(\tau_z(r)\).\\
Contrast metric \(\mathcal{C}_\tau\) compares transport efficacy across regimes.
\end{minipage}}
\caption{Schematic for Subsection~\ref{sec:ch01.6.4}: reactive and radiative regime comparison using power and torque channels.}
\label{fig:ch1.ex64-reactive-radiative}
\end{figure}

Figure~\ref{fig:ch1.ex64-reactive-radiative} clarifies why Section~1.2 regime
classification is indispensable in practice. Angular intensity or phase texture on
\(S_{r_1}\) can be large while torque transfer remains limited, because reactive
exchange dominates over exported power. At \(S_{r_2}\), conservation-compatible
transport interpretation is typically restored \cite{collin2001,jackson1999}.

This comparison assumes identical origin/axis declarations, stationary harmonic
operation, and negligible environmental scattering between \(S_{r_1}\) and
\(S_{r_2}\). Strong inhomogeneous media or multipath contamination can distort
\(P,Q,\tau_z\) trend comparison and require full propagation compensation.

The worked sequence of Section~\ref{sec:ch01.6} is therefore complete: radiation
torque conservation, phase-only counterexample, gauge-sensitivity check, and
reactive/radiative discrimination are all established in operational form.
\paragraph{Section 1.6 outcome and bridge.}
These worked problems validate the chapter's interpretive discipline under
explicit measurement geometry: torque follows conservation, phase structure alone
fails sufficiency, local gauge-sensitive maps require observability screening, and
reactive/radiative separation controls transport claims. Section~\ref{sec:ch01.7}
can now close the chapter by separating established results from deferred
developments under strict forward-only ownership.


\section{Chapter Summary and Forward Map}
\label{sec:ch01.7}

Section~\ref{sec:ch01.7} records the chapter's established outcomes, lists only
those concepts formally deferred to later owner locations, and defines the exact
inheritance boundary into Chapter~2. No new theorem content is introduced here;
the function of this section is architectural closure of Chapter~\ref{ch:01}.


\subsection{Results established in Chapter 1}
\label{sec:ch01.7.1}
% TOC tag: [DEFER]

Chapter~\ref{ch:01} closes with a fully established foundational framework for angular
electromagnetic analysis in which conservation, covariance, gauge discipline, and
measurement admissibility are no longer separate themes but a single coherent
interpretive framework. The established core begins at the Maxwell and
constitutive layer in Section~\ref{sec:ch01.1}, where local field equations,
boundary admissibility, and momentum-transfer channels were fixed as the
non-negotiable physical substrate for all later angular statements.

The chapter then defined angular structure in transport-relevant form in
Section~\ref{sec:ch01.2}. Directional content, phase/polarization topology,
transport conditions, and reactive-radiative regime separation were established in
mutual dependence rather than as isolated descriptors, thereby preventing
phase-rich patterns from being misreported as torque-capable transport when
outward-channel or conservation checks fail.

Rotational and translational rigor was established in Section~\ref{sec:ch01.3} by
fixing active-passive rotation duality, tensor/dyadic covariance, origin-shift
law, and frame-objectivity criteria. This block removes coordinate artifacts from
the chapter's admissible statement space: any claim intended for inheritance must
survive the symmetry tests formalized in that section.

Gauge structure was then resolved in Section~\ref{sec:ch01.4} by distinguishing
representation freedom from observable content. Potential transformations, local
density gauge sensitivity, dual-helicity admissibility, operational channel
definitions, and limits of local spin-orbit decomposition together established a
strict criterion: expressive decompositions are acceptable as analysis tools only
when they pass gauge/frame observability gates.

Section~\ref{sec:ch01.5} converted this theoretical framework into measurable
channel language, and Section~\ref{sec:ch01.6} tested the framework through
problem-level cases designed to expose failure modes as well as successes.
Consequently, Chapter~\ref{ch:01} now contributes a verified set of inherited
results rather than a narrative overview. Those results are complete for Volume~I
entry and are intended to be used by exact reference, not by re-derivation, in
all subsequent chapters.

\subsection{Concepts deferred to Chapters 2-4 and Volume II}
\label{sec:ch01.7.2}
% TOC tag: [DEFER]

No additional theorem is introduced in this subsection; its sole function is to
declare deferred ownership boundaries so later chapters advance without recursion.
The operator-domain and radiation-mapping formalism required for rigorous
source-to-field transformations is deferred to Chapter~2 owner locations, where
functional spaces, dyadic Green operators, spectral structure, and non-radiating
subspaces are derived once with full operator-level precision.

Representation-theoretic closure of angular eigenfields is deferred to
Chapter~3 owner locations. In particular, Hilbert/rigged-space formalization,
\(\mathrm{SO}(3)\)-linked angular momentum representation structure, and the
complete vector-spherical-harmonic eigenbasis machinery are not recast here;
Chapter~\ref{ch:01} provides only the admissibility and interpretation constraints
needed to prevent misuse when those tools are introduced.

Source realizability, finite-aperture truncation, and operator-chain limits are
deferred to Chapter~4 owner locations, where spectral truncation and physical
recoverability are treated with explicit SVD-level structure and realizability
constraints. The present chapter intentionally stops short of those derivations to
preserve one-proof ownership and avoid forward leakage of proofs.

Higher-order statistical transport, inverse-information limits, and modality
selection in complex environments are deferred beyond Volume~I by design. Their
mathematical prerequisites depend on the operator and representation layers that
Chapter~\ref{ch:01} does not own. This boundary is methodological, not editorial:
it preserves proof integrity and keeps Chapter~\ref{ch:01} as a stable foundation
for inheritance rather than a compressed survey of later theory.

\subsection{Bridge to Chapter 2}
\label{sec:ch01.7.3}
% TOC tag: [DEFER]

Chapter~2 begins from a closed foundation, not from a restart. The inheritance
contract is explicit: Maxwell-conservation structure from
Section~\ref{sec:ch01.1}, transport admissibility from
Section~\ref{sec:ch01.2}, covariance and origin discipline from
Section~\ref{sec:ch01.3}, gauge-observability boundaries from
Section~\ref{sec:ch01.4}, and operational measurement interpretation from
Sections~\ref{sec:ch01.5}--\ref{sec:ch01.6} are all carried forward by reference.

The immediate consequence is that Chapter~2 can state operator domains and
radiation mappings without re-deriving Chapter~1 laws. When Chapter~2 introduces
source-to-field operators, uniqueness classes, and Green-function actions, those
developments are interpreted through the Chapter~1 admissibility filter already
established: conservation consistency, covariance robustness, and gauge-screened
observable relevance are prerequisites, not optional commentary.

This bridge also fixes reading order for rigorous use. A Chapter~2 derivation is
accepted for onward inheritance only if its assumptions remain compatible with the
Chapter~1 boundary/validity structure. Conversely, when Chapter~2 yields new primary
results, Chapter~1 statements remain unchanged and continue to function as prior
constraints. The flow is strictly forward: Chapter~1 constrains Chapter~2
interpretation; Chapter~2 extends Chapter~1 analytic reach.

Accordingly, the transition into Chapter~2 is not thematic but structural.
Chapter~\ref{ch:01} has completed the epistemic filtering layer; Chapter~2 now
develops operator machinery on top of that layer, preserving one-proof ownership
and non-recursive progression across the remainder of Volume~I.
\paragraph{Section 1.7 terminal summary.}
Chapter~\ref{ch:01} now stands as the sole foundational basis for the core
Maxwell-conservation-gauge-observability structure used by subsequent chapters.
Chapter~2 inherits this base by reference and develops operator-domain and
radiation-mapping machinery without recursive re-derivation.


\bibliographystyle{IEEEtran}
\bibliography{references}

